%% ============================================================
%%  Scientific Watch Agent — Mémoire de fin d'études (PFE)
%%  Master Informatique / IA
%%  Compilé avec: pdflatex + biber
%%  Overleaf : xelatex ou lualatex (pour UTF-8 natif)
%% ============================================================

\documentclass[12pt,a4paper,twoside]{report}

%% ── Encodage & langue ──────────────────────────────────────
\usepackage[utf8]{inputenc}
\usepackage[T1]{fontenc}
\usepackage[french]{babel}
\usepackage{csquotes}          % requis par biblatex + french

%% ── Mise en page ───────────────────────────────────────────
\usepackage[
  top=2.5cm, bottom=2.5cm,
  left=3cm,  right=2.5cm,
  headheight=15pt
]{geometry}
\usepackage{setspace}
\onehalfspacing

%% ── En-têtes / pieds de page ───────────────────────────────
\usepackage{fancyhdr}
\pagestyle{fancy}
\fancyhf{}
\fancyhead[LE,RO]{\thepage}
\fancyhead[RE]{\small\leftmark}
\fancyhead[LO]{\small\rightmark}
\renewcommand{\headrulewidth}{0.4pt}

%% ── Typographie ────────────────────────────────────────────
\usepackage{lmodern}
\usepackage{textcomp}          % \texttimes, \textdegree, etc.
\usepackage{microtype}

%% ── Mathématiques ──────────────────────────────────────────
\usepackage{amsmath,amssymb,amsthm}
\usepackage{mathtools}

%% ── Figures & tableaux ─────────────────────────────────────
% DRAFT : remplacer par \usepackage{graphicx} quand les figures sont uploadées
\usepackage[draft]{graphicx}
\graphicspath{{figures/}}
\usepackage{float}
\usepackage{caption}
\usepackage{subcaption}
\usepackage{booktabs}
\usepackage{longtable}
\usepackage{multirow}
\usepackage{array}
\usepackage{xcolor}
\usepackage{colortbl}
\definecolor{lightgray}{gray}{0.92}
\definecolor{darkblue}{RGB}{0,51,102}

%% ── Code source ────────────────────────────────────────────
\usepackage{listings}
\lstset{
  basicstyle=\footnotesize\ttfamily,
  breaklines=true,
  frame=single,
  captionpos=b,
  language=Python,
  keywordstyle=\bfseries\color{darkblue},
  commentstyle=\itshape\color{gray},
  numbers=left, numberstyle=\tiny, numbersep=5pt,
}

%% ── Liens & PDF ────────────────────────────────────────────
\usepackage[
  colorlinks=true,
  linkcolor=darkblue,
  citecolor=darkblue,
  urlcolor=darkblue,
  pdftitle={Scientific Watch Agent : Veille Scientifique Multi-Agents avec Méta-Learning Cross-Domaine},
  pdfauthor={[Nom de l'étudiant]},
]{hyperref}

%% ── Bibliographie (IEEE) ───────────────────────────────────
\usepackage[
  backend=biber,
  style=ieee,
  sorting=none,
  maxbibnames=6,
]{biblatex}
\addbibresource{references.bib}

%% ── Algorithmes ────────────────────────────────────────────
\usepackage{algorithm}
\usepackage{algpseudocode}

%% ── Utilitaires ────────────────────────────────────────────
\usepackage{enumitem}
\usepackage{tcolorbox}
\tcbuselibrary{skins,breakable}

%% ── Définitions théorèmes ──────────────────────────────────
\newtheorem{definition}{Définition}[chapter]
\newtheorem{hypothese}{Hypothèse}[chapter]
\newtheorem{propriete}{Propriété}[chapter]

%% ── Encadrés de résultats clés ─────────────────────────────
\newtcolorbox{resultbox}{
  enhanced, breakable,
  colback=blue!5!white, colframe=darkblue,
  fonttitle=\bfseries, title=Résultat clé,
  boxrule=0.8pt, arc=3pt,
}

%% ============================================================
%%  DOCUMENT
%% ============================================================
\begin{document}

%% ── Page de titre ──────────────────────────────────────────
\begin{titlepage}
  \centering
  \vspace*{1cm}
  {\Large [Nom de l'Université / École]}\\[0.5cm]
  {\large Département Informatique}\\[2cm]

  \rule{\textwidth}{1.5pt}\\[0.4cm]
  {\LARGE\bfseries Scientific Watch Agent}\\[0.3cm]
  {\Large\bfseries Système de Veille Scientifique Multi-Agents\\
  avec Méta-Learning Cross-Domaine}\\[0.3cm]
  \rule{\textwidth}{1.5pt}\\[2cm]

  {\large Mémoire de fin d'études (PFE)\\
  Master Informatique — Spécialité Intelligence Artificielle}\\[2cm]

  {\large\bfseries [Nom Prénom]}\\[1cm]
  {\large Encadrant : [Nom de l'encadrant]}\\[2cm]

  {\large Année universitaire 2025--2026}
  \vfill
\end{titlepage}

%% ── Résumé ─────────────────────────────────────────────────
\chapter*{Résumé}
\addcontentsline{toc}{chapter}{Résumé}

Ce mémoire présente \textbf{Scientific Watch Agent}, un système multi-agents fondé
sur des grands modèles de langage (LLM) conçu pour automatiser la veille scientifique.
À partir d'un sujet de recherche, le système extrait, filtre, résume et synthétise
automatiquement la littérature académique pertinente.

La contribution centrale du mémoire est la démonstration de la \textbf{généralisation
cross-domaine} : le système s'adapte à de nouveaux domaines de recherche sans
ré-entraînement spécifique, grâce à une approche de méta-apprentissage qui prédit
les paramètres de configuration optimaux à partir de caractéristiques du domaine
(\textit{meta-features}).

Les principales contributions techniques sont :
\begin{itemize}
  \item Un pipeline multi-agents complet orchestré par LangGraph, intégrant sept
    agents spécialisés (QueryExpander, QualityCritic, Summarizer, Synthesizer, etc.) ;
  \item Un scorer de qualité scientifique à six dimensions (venue, auteurs, impact,
    vélocité, récence, pertinence) optimisé par AutoML (Optuna TPE) ;
  \item Les patterns Reflexion et ReAct appliqués à la synthèse et à l'expansion
    de requêtes ;
  \item Une évaluation méta-learning sur 19 domaines par validation croisée
    leave-one-out (LOO-CV), confirmant l'hypothèse centrale H3
    (\textit{Borda Count} : 15/19 = 79\,\%, $p = 0.014$, effet large $r_{rb} = +0.57$).
\end{itemize}

\textbf{Mots-clés :} veille scientifique, systèmes multi-agents, LLM, LangGraph,
méta-learning, AutoML, Optuna, Reflexion, ReAct, NDCG.

\chapter*{Abstract}
\addcontentsline{toc}{chapter}{Abstract}

This dissertation presents \textbf{Scientific Watch Agent}, a multi-agent system
based on large language models (LLMs) designed to automate scientific literature
monitoring. Given a research topic, the system automatically retrieves, filters,
summarises, and synthesises relevant academic publications.

The central contribution is the demonstration of \textbf{cross-domain generalisation}:
the system adapts to new research domains without domain-specific retraining,
through a meta-learning approach that predicts optimal configuration parameters
from domain characteristics (meta-features).

Key technical contributions include a seven-agent LangGraph pipeline, a six-feature
quality scorer optimised by Optuna TPE, Reflexion and ReAct patterns, and a
meta-learning evaluation across 19 domains confirming hypothesis H3
(Borda Count: 15/19 = 79\,\%, $p = 0.014$, large effect $r_{rb} = +0.57$).

\textbf{Keywords:} scientific monitoring, multi-agent systems, LLM, LangGraph,
meta-learning, AutoML, Optuna, Reflexion, ReAct, NDCG.

%% ── Remerciements ──────────────────────────────────────────
\chapter*{Remerciements}
\addcontentsline{toc}{chapter}{Remerciements}

[À compléter par l'étudiant]

%% ── Tables ─────────────────────────────────────────────────
\tableofcontents
\listoffigures
\listoftables

%% ── Chapitres ──────────────────────────────────────────────
%% ============================================================
%%  Chapitre 1 — Introduction
%% ============================================================
\chapter{Introduction}
\label{ch:introduction}

\section{Contexte et motivation}

La production scientifique mondiale connaît une croissance exponentielle :
selon les estimations, plus de deux millions d'articles de recherche sont
publiés chaque année dans les seules disciplines relevant des sciences et
technologies de l'information \cite{priem2022openalex}. Pour un chercheur
débutant sa thèse ou souhaitant explorer un domaine adjacent au sien, cette
abondance constitue paradoxalement un obstacle. Il est désormais impossible
de lire l'intégralité de la littérature pertinente ; la sélection, la
hiérarchisation et la synthèse des travaux existants sont devenues des tâches
à part entière, chronophages et susceptibles d'erreurs d'omission.

Les pratiques actuelles de veille scientifique reposent principalement sur des
approches manuelles ou semi-automatisées : interrogation des moteurs de
recherche académiques (Google Scholar, Semantic Scholar, Web of Science),
lecture d'articles de revue (\textit{survey papers}), alertes par mots-clés,
ou recommandation par les pairs. Ces approches présentent plusieurs limitations
majeures. Premièrement, elles requièrent un investissement temporel considérable,
estimé à plusieurs semaines pour une revue de littérature exhaustive dans un
nouveau domaine \cite{kitchenham2007guidelines}. Deuxièmement, elles sont
sujettes à des biais de sélection, en particulier pour les chercheurs non
initiés qui manquent des critères experts pour évaluer la qualité d'un travail.
Troisièmement, elles ne permettent pas d'actualisation automatique lorsque de
nouveaux articles paraissent.

L'émergence des grands modèles de langage (LLM, \textit{Large Language Models})
\cite{openai2023gpt4, anthropic2024claude} offre de nouvelles perspectives pour
l'automatisation de ces tâches cognitives. La capacité des LLM à comprendre,
résumer et synthétiser des textes techniques ouvre la voie à des systèmes
capables d'assister le chercheur de manière substantielle. Parallèlement,
les architectures multi-agents fondées sur des LLM \cite{wang2024survey}
permettent de décomposer des tâches complexes en sous-tâches spécialisées,
d'implémenter des mécanismes de contrôle qualité, et d'exploiter les forces
complémentaires de plusieurs modèles ou stratégies de raisonnement.

\section{Problématique}

La problématique centrale de ce mémoire peut être formulée en deux questions
imbriquées :

\begin{enumerate}
  \item \textbf{Comment concevoir un système automatisé de veille scientifique}
    qui produise des résultats de qualité comparable à celle d'un expert humain,
    en termes de pertinence des articles sélectionnés, de fidélité des résumés
    et de richesse de l'analyse des tendances ?
  \item \textbf{Comment rendre ce système généralisable à de nouveaux domaines}
    sans nécessiter une reconfiguration manuelle coûteuse pour chaque nouveau
    domaine de recherche ?
\end{enumerate}

La seconde question est au cœur de la contribution scientifique de ce travail.
Un système de veille conçu spécifiquement pour un domaine (par exemple, la
détection de fausses nouvelles) exploitera des critères de qualité spécifiques
à ce domaine : l'importance de la vélocité des citations pour les sujets
d'actualité, ou le poids accordé à la pertinence sémantique pour des champs
très techniques. Transférer naïvement ces critères à un autre domaine (par
exemple, la biologie computationnelle) peut conduire à des résultats
sous-optimaux. La généralisation constitue donc un enjeu technique et
scientifique de premier ordre.

\section{Objectifs}

Ce mémoire poursuit les objectifs suivants :

\begin{enumerate}
  \item \textbf{Concevoir et implémenter} un pipeline multi-agents complet de
    veille scientifique, intégrant la recherche d'articles, le filtrage par
    qualité, la summarisation, la synthèse et l'analyse des tendances.

  \item \textbf{Optimiser automatiquement} les paramètres de scoring documentaire
    par domaine, en utilisant des techniques d'AutoML (Optuna TPE), et démontrer
    l'amélioration par rapport aux configurations expertes fixes.

  \item \textbf{Développer une approche de méta-learning} qui prédit les
    paramètres optimaux pour un nouveau domaine à partir de caractéristiques
    mesurables de ce domaine (\textit{meta-features}), sans nécessiter
    d'optimisation préalable dans ce domaine.

  \item \textbf{Évaluer rigoureusement} la généralisation cross-domaine sur
    19 domaines de recherche couvrant à la fois les domaines de l'intelligence
    artificielle et des disciplines hors-IA (biologie computationnelle,
    traitement du langage médical, finance quantitative).
\end{enumerate}

\section{Hypothèses de recherche}
\label{sec:hypotheses}

Cinq hypothèses structurent l'évaluation expérimentale de ce mémoire.
Parmi celles-ci, deux ont été confirmées et constituent les résultats
centraux du travail :

\begin{hypothese}[H1 — AutoML par domaine]
  Les pondérations apprises par Optuna sur un domaine donné surpassent
  les pondérations expertes fixes sur ce même domaine, avec une amélioration
  relative d'au moins 5\,\% en NDCG@15.
\end{hypothese}

\begin{hypothese}[H3 — Méta-learning cross-domaine]
  La prédiction des pondérations optimales pour un nouveau domaine via
  une fonction $f_\theta : \text{meta\_features}(D) \to \text{weights}$,
  apprise sur d'autres domaines, surpasse le Transfer Direct
  (utilisation des pondérations moyennes) sur au moins 75\,\% des
  domaines de test.
\end{hypothese}

Les hypothèses H2, H4 et H5 portent respectivement sur la dégradation
prévisible du Transfer Direct, l'identification des goulots de généralisation
par l'architecture multi-agents, et la comparaison des profils de généralisation
de Claude versus GPT. Ces dernières hypothèses feront l'objet d'investigations
dans les travaux futurs (Chapitre~\ref{ch:conclusion}).

\section{Contributions}

Les contributions de ce mémoire sont les suivantes :

\paragraph{Contribution 1 — Système de veille scientifique opérationnel.}
Nous présentons un système complet, généraliste et déployable, fondé sur
sept agents LLM orchestrés par LangGraph, intégrant quatre fournisseurs
de LLM (Anthropic, OpenAI, Groq, Nebius) via une abstraction unifiée.

\paragraph{Contribution 2 — Scorer de qualité multi-dimensionnel et AutoML.}
Nous proposons un scorer à six dimensions calibré par domaine via
l'algorithme TPE d'Optuna, atteignant une amélioration moyenne de
$+68.7\,\%$ en NDCG@15 par rapport aux pondérations par défaut (H1 confirmée).

\paragraph{Contribution 3 — Méta-learning pour la généralisation cross-domaine.}
Nous montrons qu'un ensemble de méta-learners (\textit{Borda Count} combinant
kNN et BayesianRidge) prédit efficacement les pondérations optimales pour
de nouveaux domaines, surpassant le Transfer Direct sur 15 domaines sur 19
(79\,\%, test de Wilcoxon : $W^+ = 149$, $p = 0.014$, effet large $r_{rb} = +0.57$).

\paragraph{Contribution 4 — Patterns d'agents avancés.}
Nous appliquons les patterns Reflexion \cite{shinn2023reflexion} et
ReAct \cite{yao2022react} au contexte de la veille scientifique, et
évaluons l'optimisation automatique des prompts via ProTeGi
\cite{pryzant2023automatic}.

\section{Organisation du mémoire}

Ce mémoire est organisé comme suit :

\begin{itemize}
  \item Le \textbf{Chapitre~\ref{ch:etat_art}} présente l'état de l'art en
    matière de systèmes multi-agents LLM, d'AutoML pour la configuration
    de systèmes, de méta-learning, et de veille scientifique automatisée.

  \item Le \textbf{Chapitre~\ref{ch:methodologie}} décrit la méthodologie
    générale : la source de données (OpenAlex), les six dimensions de qualité,
    les oracles d'évaluation et le protocole expérimental.

  \item Le \textbf{Chapitre~\ref{ch:automl}} présente le scorer de qualité
    et son optimisation par AutoML, avec les résultats de l'hypothèse H1
    sur 19 domaines.

  \item Le \textbf{Chapitre~\ref{ch:agents}} décrit l'architecture
    multi-agents : pipeline LangGraph, contrat de communication Pydantic,
    et les sept agents spécialisés.

  \item Le \textbf{Chapitre~\ref{ch:patterns}} détaille les patterns avancés :
    ReAct (expansion de requêtes), Reflexion (révision de synthèse), et
    ProTeGi (optimisation de prompts).

  \item Le \textbf{Chapitre~\ref{ch:metalearning}} présente la contribution
    centrale : l'extraction de meta-features, les méta-learners évalués,
    et la validation de l'hypothèse H3 par test de Wilcoxon.

  \item Le dernier chapitre conclut et propose des pistes pour les
    travaux futurs, notamment la Phase 7 (agents concurrents Best-of-N)
    et le renforcement du protocole expérimental.
\end{itemize}

%% ============================================================
%%  Chapitre 2 — État de l'art
%% ============================================================
\chapter{État de l'art}
\label{ch:etat_art}

\section{Grands modèles de langage et agents LLM}

\subsection{Des LLM aux agents}

Les grands modèles de langage (LLM), tels que GPT-4 \cite{openai2023gpt4},
Claude \cite{anthropic2024claude} ou LLaMA \cite{touvron2023llama}, ont
transformé le traitement automatique du langage naturel. Ces modèles, pré-entraînés
sur de vastes corpus textuels, exhibent des capacités remarquables en matière de
compréhension, de génération et de raisonnement. Cependant, leur utilisation
isolée présente des limitations importantes : absence de mémoire persistante
entre les appels, incapacité à accéder à des sources d'information externes
en temps réel, et difficulté à décomposer des tâches complexes en sous-tâches
gérables.

La notion d'\textbf{agent LLM} adresse ces limitations en dotant le LLM de
capacités supplémentaires : accès à des outils (\textit{tool use}),
mémoire externe, et capacité à planifier une séquence d'actions. Un agent
est formellement défini comme un système capable de percevoir son environnement
et d'agir sur lui pour atteindre un objectif \cite{wooldridge1995intelligent}.
Dans le contexte des LLM, un agent correspond à un modèle de langage augmenté
d'un ensemble d'outils (API, bases de données, moteurs de recherche) et d'un
mécanisme de contrôle déterminant quand et comment les utiliser.

\subsection{Patterns de raisonnement des agents}

Plusieurs paradigmes de raisonnement ont été proposés pour structurer le
comportement des agents LLM.

\paragraph{ReAct (Reason + Act).}
Yao et al. \cite{yao2022react} proposent un cadre interleaved où l'agent
alterne entre des étapes de \textit{Thought} (raisonnement verbalisé),
d'\textit{Action} (appel à un outil) et d'\textit{Observation} (résultat
de l'outil). Ce pattern améliore significativement la précision sur des
tâches de questions-réponses nécessitant des recherches factuelles,
tout en rendant le processus de raisonnement interprétable.

\paragraph{Reflexion.}
Shinn et al. \cite{shinn2023reflexion} introduisent Reflexion, un mécanisme
par lequel un agent LLM peut s'auto-évaluer et réviser ses productions en
s'appuyant sur un signal de feedback verbal. Contrairement au fine-tuning
traditionnel qui modifie les poids du modèle, Reflexion stocke les retours
dans la mémoire épisodique et les exploite lors des itérations suivantes.
Ce pattern présente l'avantage d'améliorer la qualité des sorties sans
nécessiter d'entraînement supplémentaire.

\paragraph{Tree-of-Thought (ToT).}
Yao et al. \cite{yao2023tree} généralisent la chaîne de pensée (\textit{chain
of thought}) à une exploration arborescente de l'espace de raisonnement. L'agent
génère plusieurs chemins de pensée en parallèle, les évalue, et sélectionne le
plus prometteur. Ce paradigme excelle pour les problèmes de planification à
plusieurs étapes.

\paragraph{Self-Consistency.}
Wang et al. \cite{wang2022self} montrent que la génération de plusieurs
chemins de raisonnement indépendants suivie d'un vote majoritaire améliore
la précision des LLM sur des tâches de raisonnement arithmétique et logique.
Ce principe est à la base des approches de Best-of-N que nous envisageons
pour les travaux futurs.

\subsection{Systèmes multi-agents LLM}

Au-delà des agents individuels, les \textbf{systèmes multi-agents LLM}
(LLM-MAS) mettent en jeu plusieurs agents spécialisés qui collaborent
pour résoudre des tâches complexes. Wang et al. \cite{wang2024survey}
proposent une taxonomie distinguant les architectures coopératives
(les agents poursuivent un objectif commun) des architectures compétitives
(les agents débattent pour converger vers une réponse de meilleure qualité).

Du et al. \cite{du2023improving} démontrent que le débat multi-agents
améliore la précision factuelle des LLM sur des benchmarks de mathématiques
et de raisonnement. Les agents génèrent indépendamment leurs réponses,
puis les échangent itérativement pour raffiner leurs positions jusqu'à
convergence.

\section{Orchestration d'agents : LangGraph}

\subsection{LangChain et l'écosystème agent}

LangChain \cite{chase2022langchain} est un framework Python qui simplifie
le chaînage d'appels aux LLM et l'intégration d'outils. Il propose des
abstractions pour les chaînes (\textit{chains}), les mémoires (\textit{memory})
et les agents. Cependant, pour des pipelines complexes avec des boucles,
des conditions et un état partagé entre agents, LangChain seul s'avère
insuffisant.

\subsection{LangGraph : graphes d'état pour les agents}

LangGraph \cite{langgraph2024} est une extension de LangChain qui modélise
les pipelines agents comme des graphes orientés dont les nœuds sont des
agents et les arêtes représentent les transitions d'état. L'état du système
est un dictionnaire partagé (\texttt{TypedDict}) que chaque agent peut lire
et écrire, dans le respect de la règle \textit{Single Writer / Multiple
Readers}. Cette architecture offre trois avantages majeurs par rapport
aux approches séquentielles : la gestion naturelle des boucles (Reflexion,
ReAct), la possibilité d'arêtes conditionnelles (\textit{conditional edges})
pour router dynamiquement le flux selon l'état, et la persistance de l'état
entre les nœuds sans besoin de variables globales.

\section{Qualité des publications scientifiques}

\subsection{Indicateurs bibliométriques}

L'évaluation de la qualité d'une publication académique s'appuie
traditionnellement sur plusieurs indicateurs bibliométriques.

\paragraph{Impact de la revue.}
Le facteur d'impact (\textit{Impact Factor}) \cite{garfield1972citation}
et le classement par quartile (Q1, Q2, Q3, Q4 selon Scimago ou Web of Science)
mesurent la visibilité d'une revue dans sa discipline. Un article publié dans
une revue de rang Q1 bénéficie d'une présomption de qualité supérieure,
bien que ce critère doive être nuancé par la culture de publication propre
à chaque discipline.

\paragraph{Citations.}
Le nombre de citations d'un article constitue le proxy le plus commun de
son impact \cite{hirsch2005index}. L'indice H (\textit{h-index}) de Hirsch,
qui mesure la productivité et l'influence d'un chercheur, est un indicateur
complémentaire permettant d'évaluer l'autorité scientifique des auteurs
d'un article.

\paragraph{Vélocité des citations.}
Au-delà du compte brut, la \textit{vélocité de citation} — nombre de citations
récentes par unité de temps — permet d'identifier les travaux émergents à fort
impact potentiel, notamment les articles récents qui accumulent rapidement
des citations \cite{ke2015defining}.

\subsection{Pertinence sémantique}

La pertinence sémantique d'un article par rapport à un sujet de recherche
peut être mesurée à deux niveaux. L'approche lexicale, basée sur les
n-grammes et la correspondance de mots-clés \cite{jones1972statistical},
est rapide mais souffre du problème de synonymie et de polysémie.
L'approche par embeddings, fondée sur des modèles de langage comme
\textit{Sentence-BERT} \cite{reimers2019sentencebert} ou ses variantes
légères (all-MiniLM-L6-v2 \cite{wang2020minilm}), capture la similarité
sémantique dans un espace vectoriel continu et s'avère bien plus robuste.

\section{AutoML et optimisation d'hyperparamètres}

\subsection{Optimisation bayésienne}

L'optimisation bayésienne \cite{snoek2012practical} est une méthode de
recherche d'hyperparamètres fondée sur la construction d'un modèle de
substitution (\textit{surrogate model}) de la fonction objectif. À chaque
itération, l'algorithme sélectionne les hyperparamètres les plus prometteurs
en équilibrant exploration (zones peu connues de l'espace) et exploitation
(zones connues pour leur haute performance), via une fonction d'acquisition
telle que l'espérance d'amélioration (\textit{Expected Improvement}).

\subsection{Optuna et l'échantillonneur TPE}

Optuna \cite{akiba2019optuna} est un framework AutoML Python qui implémente
l'algorithme TPE (\textit{Tree-structured Parzen Estimator}) de Bergstra et
al.\ \cite{bergstra2011algorithms}. TPE modélise séparément la distribution
des hyperparamètres associés aux bons résultats $l(x)$ et aux mauvais résultats
$g(x)$, puis sélectionne les candidats qui maximisent le ratio $l(x)/g(x)$.
Comparé à la recherche aléatoire ou par grille, TPE converge vers de bonnes
solutions avec significativement moins d'évaluations de la fonction objectif,
ce qui est crucial lorsque chaque évaluation nécessite un appel API coûteux.

\section{Méta-learning}

\subsection{Définition et taxonomie}

Le méta-learning, ou \textit{apprentissage à apprendre}, désigne un ensemble
de techniques visant à améliorer les performances d'un algorithme d'apprentissage
en exploitant l'expérience accumulée sur de multiples tâches passées
\cite{vanschoren2018meta}. Contrairement à l'apprentissage conventionnel qui
optimise les performances sur une tâche spécifique, le méta-learning cherche à
apprendre une représentation ou une procédure générale transférable à de
nouvelles tâches avec peu d'exemples.

Trois grandes familles d'approches existent :
\begin{itemize}
  \item Les méthodes fondées sur l'optimisation du modèle initial (MAML
    \cite{finn2017model} et ses variantes), qui apprennent une initialisation
    des poids du modèle facilitant le fine-tuning rapide sur de nouvelles tâches ;
  \item Les méthodes fondées sur des métriques (Prototypical Networks
    \cite{snell2017prototypical}), qui apprennent un espace d'embedding dans
    lequel la classification par plus proche voisin est efficace ;
  \item Les méthodes fondées sur des modèles (\textit{model-based}), qui apprennent
    un mécanisme d'inférence rapide, souvent sous la forme d'un réseau de
    mémoire ou d'un transformeur.
\end{itemize}

\subsection{Méta-learning pour la configuration de systèmes}

Notre travail s'inscrit dans un paradigme de méta-learning différent des
approches précédentes, plus proche du domaine AutoML \cite{hutter2019automated}
et du \textit{meta-learning for algorithm configuration}
\cite{feurer2015initializing}. L'idée est de prédire les hyperparamètres
optimaux d'un algorithme (ici, les pondérations du scorer) pour une nouvelle
tâche (ici, un nouveau domaine de recherche), à partir de caractéristiques
descriptives de cette tâche (\textit{meta-features}).

Feurer et al.\ \cite{feurer2015initializing} utilisent les meta-features
d'un dataset (nombre d'instances, de features, ratio de classes, etc.) pour
initialiser intelligemment la recherche d'hyperparamètres dans Auto-sklearn.
Dans notre contexte, les meta-features décrivent le corpus de publications
d'un domaine (distribution des citations, niveau des auteurs, récence des
articles, etc.) et permettent de prédire les pondérations optimales du scorer.

\subsection{Combinaison d'estimateurs : agrégation de rangs}

La combinaison de plusieurs estimateurs (\textit{ensemble methods}) améliore
typiquement la robustesse et la performance par rapport à un estimateur unique.
Dans le contexte de la prédiction de vecteurs de pondérations, nous adaptons
des méthodes d'agrégation de rangs initialement développées pour la fusion
de résultats de recherche d'information.

Le \textbf{comptage de Borda} \cite{borda1784memoire}, formalisé dans un
contexte d'apprentissage par Dwork et al.\ \cite{dwork2001rank}, attribue
à chaque alternative un score égal à sa position dans le classement inverse,
puis somme ces scores sur tous les votants. Il s'agit d'une méthode robuste
aux outliers et naturellement insensible aux différences d'échelle entre
estimateurs.

La \textbf{fusion par rang réciproque} (RRF, \textit{Reciprocal Rank Fusion})
\cite{cormack2009reciprocal} assigne à chaque alternative un score
$\text{RRF}(d) = \sum_{r \in R} \frac{1}{k + \text{rank}_r(d)}$ où $k$ est
un paramètre de lissage. RRF est particulièrement efficace pour la fusion
de listes de rangs provenant de systèmes hétérogènes.

\section{Métriques d'évaluation en recherche d'information}

\subsection{NDCG (\textit{Normalized Discounted Cumulative Gain})}

La métrique NDCG \cite{jarvelin2002cumulated} est le standard d'évaluation
des systèmes de classement en recherche d'information. Elle prend en compte
à la fois la pertinence et la position des documents dans le classement,
en appliquant un facteur de discount logarithmique aux positions inférieures.

Pour une liste de $k$ documents, le DCG est défini par :
\begin{equation}
  \text{DCG}@k = \sum_{i=1}^{k} \frac{2^{r_i} - 1}{\log_2(i+1)}
  \label{eq:dcg}
\end{equation}
où $r_i \in \{0, 1, 2\}$ est la pertinence du document en position $i$
(0 = non pertinent, 1 = pertinent, 2 = très pertinent).

Le NDCG est obtenu en normalisant le DCG par le DCG idéal (IDCG) :
\begin{equation}
  \text{NDCG}@k = \frac{\text{DCG}@k}{\text{IDCG}@k}
  \label{eq:ndcg}
\end{equation}

Nous utilisons NDCG@15 comme métrique principale, ce choix étant motivé
par la pratique usuelle de consulter les 15 à 20 premiers résultats d'une
recherche académique.

\subsection{Tests statistiques non-paramétriques}
\label{sec:etat_art_wilcoxon}

Le \textbf{test de Wilcoxon signé} \cite{wilcoxon1945individual} est une
alternative non-paramétrique au test $t$ apparié, adaptée aux situations
où les hypothèses de normalité ne peuvent être vérifiées. Pour $n$ paires
d'observations $(x_i, y_i)$, on considère les différences $d_i = y_i - x_i$,
on les classe par valeur absolue, et on calcule la statistique $W^+$ (somme
des rangs des différences positives). La corrélation rang-bisériale
$r_{rb} = (W^+ - W^-) / (W^+ + W^-)$ mesure la taille d'effet, avec les
conventions : $|r_{rb}| < 0.3$ (petit), $0.3$--$0.5$ (moyen), $> 0.5$ (grand).

\section{Optimisation automatique des prompts}

\subsection{L'importance du \textit{prompt engineering}}

La qualité des sorties d'un LLM dépend fortement de la formulation des
instructions (\textit{prompts}) qui lui sont soumises. Le \textit{prompt
engineering} manuel est un processus itératif et subjectif, difficile à
systématiser. Plusieurs travaux ont proposé d'automatiser cette optimisation.

\subsection{ProTeGi}

Pryzant et al.\ \cite{pryzant2023automatic} proposent ProTeGi
(\textit{Prompt Optimization with Textual Gradients}), une méthode qui
adapte le concept de gradient de descente au domaine textuel. ProTeGi
utilise un LLM ``critique'' pour identifier les faiblesses d'un prompt
et proposer des modifications ciblées, de façon analogue à un gradient
qui indique la direction d'amélioration. Cette approche ne nécessite
pas d'accès aux gradients du modèle et est donc applicable à tout LLM
accessible via API.

\section{Veille scientifique automatisée — travaux existants}

\subsection{Systèmes de recommandation académique}

Plusieurs systèmes automatisés de veille ou de recommandation académique
ont été proposés. \textit{Semantic Scholar} \cite{lo2020s2orc} utilise des
embeddings pour recommander des articles similaires. \textit{ResearchRabbit}
et \textit{Connected Papers} proposent des visualisations de graphes de
citations pour explorer des corpus. Ces systèmes se concentrent sur la
découverte d'articles similaires à un corpus de départ, sans aborder
l'analyse automatique du contenu ni la synthèse.

\subsection{Génération automatique de revues de littérature}

Des travaux récents ont exploré la génération automatique de revues de
littérature (\textit{automated survey generation}) à l'aide de LLM.
Liu et al.\ \cite{liu2023survey} utilisent un agent LLM pour décomposer
un sujet de recherche en sous-thèmes et récupérer les articles pertinents
pour chacun. Cependant, ces approches ne proposent pas de mécanisme de
filtrage par qualité ni d'évaluation rigoureuse de la généralisation à de
nouveaux domaines.

\subsection{Positionnement de notre travail}

Par rapport aux travaux existants, notre contribution se distingue sur
trois axes :
\begin{enumerate}
  \item \textbf{Scoring multi-dimensionnel et calibration par AutoML} : nous
    formalisons l'évaluation de la qualité d'un article selon six dimensions
    et apprenons automatiquement les pondérations optimales par domaine.
  \item \textbf{Architecture multi-agents complète} : nous proposons un pipeline
    end-to-end intégrant sept agents spécialisés, des patterns de raisonnement
    avancés (ReAct, Reflexion) et un contrat de communication formalisé.
  \item \textbf{Méta-learning pour la généralisation} : nous évaluons
    rigoureusement, sur 19 domaines, la capacité du système à s'adapter
    à de nouveaux domaines sans reconfiguration manuelle.
\end{enumerate}

Le tableau~\ref{tab:related_work} synthétise le positionnement de notre
approche par rapport aux systèmes existants.

\begin{table}[htbp]
  \centering
  \caption{Comparaison des approches de veille scientifique automatisée.}
  \label{tab:related_work}
  \rowcolors{2}{lightgray}{white}
  \begin{tabular}{lcccc}
    \toprule
    \textbf{Système} & \textbf{Multi-agents} & \textbf{Scoring} & \textbf{Synthèse} & \textbf{Généralisation} \\
    \midrule
    Semantic Scholar & \texttimes & partiel & \texttimes & \checkmark \\
    Liu et al. (2023) & \checkmark & \texttimes & \checkmark & \texttimes \\
    ResearchRabbit & \texttimes & \texttimes & \texttimes & \checkmark \\
    \textbf{Notre approche} & \checkmark & \checkmark & \checkmark & \checkmark \\
    \bottomrule
  \end{tabular}
\end{table}

%% ============================================================
%%  Chapitre 3 — Méthodologie générale
%% ============================================================
\chapter{Méthodologie générale}
\label{ch:methodologie}

\section{Vue d'ensemble du système}

\textit{Scientific Watch Agent} s'articule autour d'un pipeline en cinq
étapes principales, illustré à la figure~\ref{fig:pipeline_overview} :

\begin{enumerate}
  \item \textbf{Recherche} (\textit{Search}) : expansion de la requête et
    récupération d'articles depuis OpenAlex ;
  \item \textbf{Filtrage par qualité} (\textit{Quality Filter}) : scoring
    multi-dimensionnel et sélection des articles les plus pertinents ;
  \item \textbf{Summarisation} (\textit{Summarize}) : génération de résumés
    structurés ou narratifs pour chaque article ;
  \item \textbf{Synthèse} (\textit{Synthesize}) : génération d'une vue
    d'ensemble du domaine à partir des résumés ;
  \item \textbf{Analyse des tendances} (\textit{Trend Analysis}) :
    identification des directions émergentes, des lacunes et des
    perspectives de recherche.
\end{enumerate}

\begin{figure}[htbp]
  \centering
  \includegraphics[width=\textwidth]{figures/pipeline_overview.pdf}
  \caption{Vue d'ensemble du pipeline Scientific Watch Agent. Chaque étape
  correspond à un agent spécialisé dans l'architecture LangGraph.}
  \label{fig:pipeline_overview}
\end{figure}

\section{Source de données : OpenAlex}

\subsection{Présentation d'OpenAlex}

OpenAlex \cite{priem2022openalex} est un index open-source de la littérature
académique mondiale, développé par l'organisation non gouvernementale OurResearch.
Il indexe plus de 250 millions de publications, d'auteurs, d'institutions
et de sources, en exploitant des données publiques telles que Crossref,
PubMed, ORCID et Microsoft Academic Graph.

OpenAlex présente plusieurs avantages déterminants pour notre usage.
Premièrement, son accès est entièrement gratuit via une API REST, contrairement
à des alternatives comme Web of Science ou Scopus qui nécessitent des abonnements
institutionnels onéreux. Deuxièmement, il offre des données structurées riches :
informations sur les auteurs (h-index, affiliations), les venues (classement
par quartile), les citations entrantes et sortantes. Troisièmement, il fournit
une API de filtre avancée permettant des requêtes par concept, domaine,
plage d'années et nombre minimum de citations.

\subsection{Interface OpenAlex dans le système}

Notre module \texttt{src/sources/openalex.py} implémente un client Python
qui encapsule les appels à l'API OpenAlex. Il gère notamment le \textit{cursor
pagination} pour récupérer de larges corpus (jusqu'à plusieurs centaines
d'articles par requête), la déduplication par identifiant OpenAlex,
et le mapping vers le schéma \texttt{Article} défini dans \texttt{src/schemas.py}.

La recherche est effectuée par concepts (\texttt{/concepts}),
en filtrant sur la plage d'années et en triant par nombre de citations
décroissant. Pour chaque article récupéré, les informations d'auteur
(h-index, institution) sont enrichies via des appels secondaires à
\texttt{/authors/\{id\}}, avec une gestion de cache LRU pour limiter
le coût des appels réseau.

\section{Dimensions de qualité : le scorer à six features}
\label{sec:scoring_features}

Nous évaluons la qualité d'un article selon six dimensions complémentaires.
Le score final est une combinaison linéaire pondérée :

\begin{equation}
  \text{score}(a) = \sum_{f \in \mathcal{F}} w_f \cdot s_f(a)
  \label{eq:weighted_score}
\end{equation}

où $\mathcal{F} = \{\text{venue}, \text{authors}, \text{impact},
\text{velocity}, \text{recency}, \text{relevance}\}$,
$s_f(a) \in [0, 1]$ est le sous-score normalisé pour la feature $f$,
et $w_f \geq 0$ avec $\sum_f w_f = 1$.

\subsection{Score de venue (\texttt{venue\_score})}

Le score de venue mesure la qualité de la revue ou conférence dans laquelle
l'article a été publié, selon son classement par quartile Scimago.
Nous définissons un mapping discret : Q1 $\to$ 1.0, Q2 $\to$ 0.7,
Q3 $\to$ 0.4, Q4 $\to$ 0.2, non classé $\to$ 0.1.

Les conférences de premier plan sans classement Scimago (par exemple,
NeurIPS, ICML, CVPR) reçoivent un bonus spécifique via une liste
de venues reconnues maintenue dans \texttt{src/config.py}.

\subsection{Score d'auteurs (\texttt{authors\_score})}

Le score d'auteurs reflète l'autorité scientifique des contributeurs
de l'article, mesurée par l'indice H maximal parmi les auteurs :

\begin{equation}
  s_{\text{authors}}(a) = \frac{\log(1 + h_{\max}(a))}{\log(1 + H_{\text{ref}})}
  \label{eq:authors_score}
\end{equation}

où $h_{\max}(a)$ est l'indice H de l'auteur le plus cité parmi les
contributeurs de $a$, et $H_{\text{ref}} = 50$ est une valeur de référence
normalisatrice. L'application du logarithme atténue l'effet des valeurs
extrêmes (chercheurs à très fort indice H) et évite que ce facteur
n'écrase les autres dimensions.

\subsection{Score d'impact (\texttt{impact\_score})}

Le score d'impact mesure l'influence de l'article via son nombre de
citations rapporté à son âge :

\begin{equation}
  s_{\text{impact}}(a) = \frac{\log(1 + c(a)/\max(1, \text{age}(a)))}
    {\log(1 + C_{\text{ref}})}
  \label{eq:impact_score}
\end{equation}

où $c(a)$ est le nombre de citations, $\text{age}(a)$ est l'ancienneté
de l'article en années, et $C_{\text{ref}} = 20$ est la valeur de référence.

\subsection{Score de vélocité (\texttt{velocity\_score})}

La vélocité de citation mesure la dynamique récente de l'article, c'est-à-dire
la proportion de ses citations accumulées au cours des deux dernières années
par rapport à l'ensemble de ses citations :

\begin{equation}
  s_{\text{velocity}}(a) = \min\left(1, \frac{c_{\text{recent}}(a)}{V_{\text{ref}}}\right)
  \label{eq:velocity_score}
\end{equation}

où $c_{\text{recent}}(a)$ est le nombre de citations des deux dernières années
et $V_{\text{ref}} = 20$ est le seuil de saturation. Cette feature est
particulièrement discriminante pour les domaines d'actualité où les articles
récents à fort impact sont préférables aux travaux fondateurs anciens.

\subsection{Score de récence (\texttt{recency\_score})}

Le score de récence favorise les travaux récents, en appliquant une décroissance
exponentielle en fonction de l'âge de l'article :

\begin{equation}
  s_{\text{recency}}(a) = \exp\left(-\lambda \cdot \text{age}(a)\right)
  \label{eq:recency_score}
\end{equation}

avec une demi-vie de trois ans ($\lambda = \ln(2)/3 \approx 0.231$),
signifiant qu'un article vieux de trois ans voit son score de récence
réduit de moitié par rapport à un article de l'année courante.

\subsection{Score de pertinence (\texttt{relevance\_score})}

Le score de pertinence évalue l'adéquation thématique de l'article
avec le sujet de recherche soumis par l'utilisateur.

\paragraph{V1 — approche lexicale.}
La version lexicale (V1) calcule la correspondance entre les n-grammes
du titre et de l'abstract de l'article et les mots-clés du topic de
recherche, avec un bonus pour les bigrammes.

\paragraph{V2 — approche sémantique.}
La version sémantique (V2) encode le titre et l'abstract de l'article
ainsi que le topic dans un espace vectoriel de dimension 384 via
le modèle \textit{all-MiniLM-L6-v2} \cite{wang2020minilm}, et calcule
la similarité cosinus entre les vecteurs obtenus :

\begin{equation}
  s_{\text{relevance}}^{V2}(a) = \frac{e(a) \cdot e(q)}{\|e(a)\| \cdot \|e(q)\|}
  \label{eq:relevance_v2}
\end{equation}

où $e(a)$ est l'embedding de $a$ et $e(q)$ celui du topic $q$.
Le modèle est chargé localement (sans appel API), mis en cache via
un cache LRU, et son empreinte mémoire est de l'ordre de 80 Mo.

\section{Protocole d'évaluation}

\subsection{Construction des oracles}
\label{sec:oracle_construction}

L'évaluation d'un système de classement de publications nécessite un
oracle de pertinence : une annotation de référence indiquant, pour un
ensemble d'articles, leur niveau de pertinence vis-à-vis du sujet.

\paragraph{Principe de l'oracle.}
Pour chaque domaine $d$, nous récupérons un ensemble de 3 à 5 articles
de revue (\textit{survey papers}) via OpenAlex. Les articles cités dans
la bibliographie d'au moins deux de ces surveys sont considérés comme
\textbf{grade-2} (très pertinents, $r_i = 2$), ceux cités par un seul
survey comme \textbf{grade-1} (pertinents, $r_i = 1$). Les autres articles
du corpus récupéré pour le domaine constituent le bruit de fond
(\textbf{grade-0}, $r_i = 0$).

Cette heuristique de construction est fondée sur l'hypothèse que les
articles régulièrement cités par plusieurs surveys indépendants sont
les références incontournables du domaine — c'est-à-dire les articles
qu'un système de veille idéal devrait placer en tête de classement.

\paragraph{Réparation des oracles (\textit{oracle repair}).}
Certains domaines, notamment les plus récents ou les moins couverts par
des surveys, produisent des oracles avec peu ou pas d'articles grade-2.
Un oracle sans grade-2 rend le calcul de NDCG@15 non informatif (score
identique pour tout classement). Nous avons développé l'utilitaire
\texttt{promote\_oracle\_grade2.py} qui permet, dans ce cas, de promouvoir
manuellement des articles grade-1 sélectionnés par un expert en grade-2,
sur la base de critères de qualité explicites (articles de venues
prestigieuses, fortement cités, directement liés au topic).

\paragraph{Corpus d'évaluation.}
Au total, 19 oracles ont été construits, couvrant 16 domaines d'IA et
3 domaines hors-IA (biologie computationnelle, NLP médical, finance
quantitative). La taille des corpus varie de 90 à 350 articles, avec
un nombre de grade-2 allant de 5 à 21 selon le domaine.

\subsection{Métrique principale : NDCG@15}

Le NDCG@15 (défini à l'équation~\ref{eq:ndcg}) est choisi comme
métrique principale pour les raisons suivantes :

\begin{itemize}
  \item Il prend en compte la \textit{position} dans le classement : placer
    un article très pertinent en premier position est préférable à le
    placer en quinzième ;
  \item Il est normalisé dans $[0, 1]$, facilitant la comparaison entre
    domaines de tailles différentes ;
  \item Il supporte plusieurs niveaux de pertinence (grade-0, 1, 2),
    contrairement aux métriques binaires comme la précision ou le rappel ;
  \item Le seuil $k=15$ correspond à la pratique usuelle de consultation
    des premiers résultats d'une recherche académique.
\end{itemize}

\subsection{Protocole de validation : Leave-One-Out}

Pour l'évaluation du méta-learner (Chapitre~\ref{ch:metalearning}),
nous adoptons une validation croisée \textit{Leave-One-Out} (LOO-CV).
À chaque itération $i \in \{1, \ldots, 19\}$, le domaine $d_i$ est
exclu de l'ensemble d'entraînement, le méta-learner est ajusté sur les
18 domaines restants, et sa performance est évaluée sur $d_i$. Cette
procédure garantit qu'aucune information du domaine de test ne fuite
dans l'entraînement, et produit une estimation non biaisée de la
capacité de généralisation.

\section{Implémentation et choix technologiques}

\subsection{Contrat de communication — Pydantic v2}

Toutes les données structurées échangées entre les agents sont définies
comme des modèles Pydantic v2 \cite{pydantic2024} dans \texttt{src/schemas.py}.
Cette approche présente deux avantages : la validation automatique des types
à chaque transition entre agents, et une documentation implicite du
contrat d'interface qui facilite l'évolution indépendante des agents.

Les schémas principaux sont : \texttt{Article}, \texttt{QualityScore},
\texttt{ArticleSummary}, \texttt{NarrativeSummary}, \texttt{Synthesis},
\texttt{TrendAnalysis}, \texttt{CriticFeedback}, et \texttt{AgentLog}.

\subsection{Gestion des coûts API}

Chaque appel LLM est tracé via \texttt{AgentLog} qui enregistre le nombre
de tokens consommés, le coût en USD et la durée d'exécution. Les optimisations
de coût mises en place incluent : l'utilisation de modèles légers pour
les tâches répétitives (Nebius Qwen3-30B à \$0.14/M tokens pour la
summarisation), le prompt caching Anthropic (90\,\% de réduction sur
les system prompts récurrents), et l'utilisation de Groq pour les tâches
latence-sensibles (inférence Llama à moins de 500\,ms).

%% ============================================================
%%  Chapitre 4 — AutoML pour le scoring scientifique
%% ============================================================
\chapter{AutoML pour le scoring scientifique}
\label{ch:automl}

\section{Motivation : les limites des pondérations expertes}

La combinaison linéaire des six dimensions de qualité définie à
l'équation~\ref{eq:weighted_score} soulève une question fondamentale :
comment choisir les pondérations $w_f$ ? Une approche naïve consiste
à fixer des pondérations expertes uniformes ($w_f = 1/6 \approx 0.167$)
ou à les calibrer manuellement en se fondant sur l'intuition du domaine.

Cette approche présente deux limitations majeures. D'une part, elle suppose
que les six dimensions contribuent de manière égale à la qualité perçue,
ce qui est inexact : l'importance relative de la vélocité (pour les domaines
d'actualité) versus la récence (pour les domaines plus matures) varie
significativement selon le contexte. D'autre part, une calibration manuelle
est subjective et difficilement reproductible d'un domaine à l'autre.

L'optimisation automatique des pondérations par un algorithme AutoML
constitue une alternative rigoureuse, reproductible et potentiellement
plus performante. C'est l'objet de ce chapitre.

\section{Formalisation du problème d'optimisation}

\subsection{Espace de recherche}

L'espace de recherche est le simplexe unitaire en dimension 6 :

\begin{equation}
  \mathcal{W} = \left\{ \mathbf{w} \in \mathbb{R}^6 \mid
    w_f \geq 0 \;\forall f,\; \sum_{f} w_f = 1 \right\}
  \label{eq:simplex}
\end{equation}

En pratique, Optuna optimise les six pondérations dans $[0, 1]$
indépendamment, puis normalise le vecteur résultant pour satisfaire
la contrainte de somme unitaire. Ce reparamétrisation préserve
la topologie de l'espace de recherche tout en simplifiant l'exploration.

\subsection{Fonction objectif}

La fonction objectif est le NDCG@15 calculé sur l'oracle du domaine :

\begin{equation}
  \mathbf{w}^* = \arg\max_{\mathbf{w} \in \mathcal{W}}
    \text{NDCG@15}\!\left(\text{rank}_{\mathbf{w}}(C_d), G_d\right)
  \label{eq:optimization_objective}
\end{equation}

où $C_d$ est le corpus du domaine $d$, $G_d$ est l'oracle de pertinence
(grades 0, 1, 2 pour chaque article), et $\text{rank}_{\mathbf{w}}(C_d)$
est le classement des articles du corpus selon le score
$\text{score}_{\mathbf{w}}(a) = \sum_f w_f \cdot s_f(a)$.

\subsection{Algorithme : TPE d'Optuna}

Nous utilisons l'échantillonneur TPE (\textit{Tree-structured Parzen
Estimator}) d'Optuna \cite{akiba2019optuna}, configuré avec 150 essais
(\textit{trials}) par domaine. TPE maintient deux distributions
de probabilité : $l(\mathbf{w})$ modélisant les pondérations associées
aux meilleurs essais, et $g(\mathbf{w})$ celles associées aux essais
moins performants. À chaque itération, il sélectionne les pondérations
qui maximisent le ratio $\text{EI}(\mathbf{w}) \propto l(\mathbf{w})/g(\mathbf{w})$,
focalisant progressivement l'exploration sur les régions prometteuses
de $\mathcal{W}$.

\section{Résultats : validation de H1}

\subsection{Amélioration par domaine}

Le tableau~\ref{tab:automl_results} présente les résultats de l'optimisation
AutoML pour les 19 domaines d'évaluation. Pour chaque domaine, nous
reportons le NDCG@15 obtenu avec les pondérations par défaut (uniformes),
les pondérations apprises par Optuna (V2, avec pertinence sémantique),
et l'amélioration relative.

\begin{table}[htbp]
  \centering
  \caption{Résultats AutoML par domaine : NDCG@15 et amélioration relative.}
  \label{tab:automl_results}
  \rowcolors{2}{lightgray}{white}
  \begin{tabular}{lccr}
    \toprule
    \textbf{Domaine} & \textbf{Default} & \textbf{Optuna V2} & \textbf{Amélioration} \\
    \midrule
    fake\_news\_detection  & 0.093 & 0.470 & $+406\,\%$ \\
    neural\_architecture\_search & 0.359 & 0.590 & $+64\,\%$ \\
    knowledge\_graph\_embedding & 0.332 & 0.602 & $+81\,\%$ \\
    graph\_attention\_networks & 0.396 & 0.645 & $+63\,\%$ \\
    explainable\_ai & 0.067 & 0.206 & $+208\,\%$ \\
    text\_summarization & 0.391 & 0.681 & $+74\,\%$ \\
    computational\_biology & 0.470 & 0.566 & $+20\,\%$ \\
    sentiment\_analysis & 0.151 & 0.351 & $+133\,\%$ \\
    deep\_reinforcement\_learning & 0.124 & 0.251 & $+102\,\%$ \\
    \textit{Moyenne (19 domaines)} & \textit{0.229} & \textit{0.387} & \textit{$+68.7\,\%$} \\
    \bottomrule
  \end{tabular}
  \small{\textit{Note : seuls 9 domaines représentatifs sont listés par souci de lisibilité.
  Les 19 résultats complets sont disponibles dans \texttt{data/metalearning/phase8c\_results.json}.}}
\end{table}

\begin{resultbox}
  \textbf{H1 confirmée.} Sur les 19 domaines évalués, l'optimisation Optuna
  améliore systématiquement le NDCG@15 par rapport aux pondérations par défaut,
  avec une amélioration moyenne de $+68.7\,\%$ ($\text{min} = +20\,\%$,
  $\text{max} = +406\,\%$). L'amélioration minimale dépasse largement le
  seuil de $5\,\%$ fixé par H1, confirmant l'hypothèse.
\end{resultbox}

\subsection{Analyse des profils de pondération}

Au-delà des performances agrégées, l'analyse des vecteurs de pondérations
appris révèle des patterns intéressants qui constituent le signal brut
pour l'apprentissage méta (Chapitre~\ref{ch:metalearning}).

Le tableau~\ref{tab:weight_profiles} illustre la diversité des profils
appris sur quelques domaines représentatifs.

\begin{table}[htbp]
  \centering
  \caption{Profils de pondération appris par Optuna (pondérations normalisées).}
  \label{tab:weight_profiles}
  \rowcolors{2}{lightgray}{white}
  \begin{tabular}{lcccccc}
    \toprule
    \textbf{Domaine} & \textbf{venue} & \textbf{authors} & \textbf{impact}
      & \textbf{velocity} & \textbf{recency} & \textbf{relevance} \\
    \midrule
    fake\_news & 0.027 & 0.061 & 0.037 & \textbf{0.698} & 0.104 & 0.024 \\
    federated\_learning & 0.201 & \textbf{0.295} & 0.175 & 0.138 & 0.153 & 0.038 \\
    graph\_NN & 0.204 & 0.128 & 0.084 & 0.221 & \textbf{0.329} & 0.034 \\
    explainable\_ai & 0.089 & 0.074 & 0.214 & \textbf{0.296} & 0.029 & \textbf{0.298} \\
    comp. biology & 0.112 & \textbf{0.281} & 0.168 & 0.165 & 0.210 & 0.063 \\
    quant. finance & \textbf{0.283} & 0.150 & 0.069 & 0.114 & 0.032 & \textbf{0.352} \\
    \bottomrule
  \end{tabular}
\end{table}

Plusieurs observations se dégagent. La pondération de la \textbf{vélocité}
est prépondérante pour la détection de fausses nouvelles (0.698), un domaine
d'actualité où les articles récemment cités sont les plus pertinents.
À l'inverse, pour les \textbf{réseaux de neurones sur graphes} (domaine
plus mature), la récence prend le dessus (0.329). La \textbf{pertinence
sémantique} est fortement valorisée pour l'IA explicable et la finance
quantitative, deux domaines au vocabulaire technique très spécifique.

Ces divergences de profils constituent la manifestation concrète
du phénomène que le méta-learning cherche à modéliser :
\textit{des domaines aux caractéristiques différentes nécessitent des
stratégies de scoring différentes}.

\section{V1 vs V2 : impact de la pertinence sémantique}

L'un des choix techniques importants de ce travail est d'utiliser la
pertinence sémantique (V2, embeddings) plutôt que la pertinence lexicale
(V1, mots-clés) dans le scorer. Le tableau~\ref{tab:v1_v2} compare
les deux approches sur plusieurs domaines représentatifs.

\begin{table}[htbp]
  \centering
  \caption{Comparaison V1 (lexicale) vs V2 (sémantique) en NDCG@15.}
  \label{tab:v1_v2}
  \rowcolors{2}{lightgray}{white}
  \begin{tabular}{lccc}
    \toprule
    \textbf{Domaine} & \textbf{Optuna V1} & \textbf{Optuna V2} & \textbf{$\Delta$(V2-V1)} \\
    \midrule
    knowledge\_graph\_embedding & 0.599 & 0.602 & $+0.003$ \\
    explainable\_ai & 0.206 & 0.206 & $\pm 0.000$ \\
    computational\_biology & 0.547 & 0.566 & $+0.020$ \\
    text\_summarization & 0.263 & 0.681 & $+0.418$ \\
    quantitative\_finance & 0.099 & 0.226 & $+0.127$ \\
    \bottomrule
  \end{tabular}
\end{table}

La pertinence sémantique apporte un gain significatif pour les domaines
à vocabulaire spécialisé (biologie computationnelle, finance quantitative,
text summarization), où la correspondance lexicale stricte est insuffisante.
Pour les domaines à terminologie plus standard (knowledge graph, XAI),
les deux approches sont comparables. Ces résultats justifient le choix
de V2 comme version par défaut dans le système.

\section{Implémentation}

\subsection{Module \texttt{automl\_scorer.py}}

Le scorer AutoML est implémenté dans \texttt{src/scoring/automl\_scorer.py}.
La fonction principale \texttt{run\_optuna\_study()} orchestre :
\begin{enumerate}
  \item Le chargement de l'oracle depuis \texttt{data/oracle/\{domain\}/} ;
  \item La création d'une étude Optuna avec échantillonneur TPE et
    pruner MedianPruner ;
  \item L'optimisation de la fonction objectif NDCG@15 sur 150 essais ;
  \item La sauvegarde des pondérations optimales dans
    \texttt{data/weights/\{topic\}\_weights.json}.
\end{enumerate}

Les poids sont chargés automatiquement par \texttt{QualityCritic}
lors de l'exécution du pipeline, via la fonction
\texttt{load\_weights\_for\_topic(topic)}.

\subsection{Accélération par sous-scores précalculés}

Afin d'accélérer l'optimisation, les sous-scores ($s_f(a)$ pour chaque
article $a$ et feature $f$) sont calculés une seule fois avant la boucle
d'optimisation et mis en cache. L'évaluation d'un vecteur de pondérations
$\mathbf{w}$ se réduit alors à une multiplication matrice-vecteur suivie
du calcul de NDCG@15, ce qui prend moins d'une milliseconde. Avec 150
essais, l'optimisation complète d'un domaine nécessite moins de 30 secondes.

%% ============================================================
%%  Chapitre 5 — Architecture multi-agents
%% ============================================================
\chapter{Architecture multi-agents}
\label{ch:agents}

\section{Principes de conception}

L'architecture multi-agents de \textit{Scientific Watch Agent} repose sur
quatre principes directeurs.

\paragraph{Séparation des responsabilités.}
Chaque agent a une responsabilité unique et bien délimitée : recherche,
filtrage, résumé, synthèse, critique, analyse. Cette séparation facilite
le test, la maintenance et l'évolution indépendante de chaque composant.

\paragraph{Contrat de communication formalisé.}
Tous les échanges entre agents transitent par un état partagé (\texttt{WatchState})
dont chaque champ est typé par un schéma Pydantic v2. Ce contrat formalisé
garantit la cohérence des données tout au long du pipeline et constitue
une documentation vivante de l'interface entre agents.

\paragraph{Principe Single Writer / Multiple Readers.}
Chaque section de l'état partagé n'est écrite que par un unique agent,
mais peut être lue par tous. Cette règle prévient les conditions de course
(\textit{race conditions}) dans les architectures futures à agents concurrents.

\paragraph{Abstraction du fournisseur LLM.}
Les agents ne connaissent pas le fournisseur de LLM qu'ils utilisent.
Ils invoquent une interface unifiée \texttt{LLMClient} qui délègue
automatiquement au fournisseur approprié selon la tâche, via une fabrique
(\texttt{factory.py}).

\section{État partagé : \texttt{WatchState}}

L'état partagé \texttt{WatchState} est défini comme un \texttt{TypedDict}
Python dans \texttt{src/agents/state.py}. Il contient :

\begin{itemize}
  \item \texttt{topic} : le sujet de recherche soumis par l'utilisateur ;
  \item \texttt{expanded\_queries} : liste des requêtes générées par
    le \texttt{QueryExpander} ;
  \item \texttt{raw\_articles} : liste d'objets \texttt{Article} récupérés
    par le \texttt{Searcher} ;
  \item \texttt{scored\_articles} : liste d'articles avec leur
    \texttt{QualityScore} associé, produite par le \texttt{QualityCritic} ;
  \item \texttt{summaries} : liste d'\texttt{ArticleSummary} ou
    \texttt{NarrativeSummary} générés par le \texttt{Summarizer} ;
  \item \texttt{synthesis} : objet \texttt{Synthesis} produit par
    le \texttt{Synthesizer} ;
  \item \texttt{critic\_feedback} : objet \texttt{CriticFeedback}
    du dernier cycle Reflexion ;
  \item \texttt{reflexion\_count} : compteur d'itérations Reflexion ;
  \item \texttt{trend\_analysis} : objet \texttt{TrendAnalysis} final ;
  \item \texttt{logs}, \texttt{errors} : listes accumulées via
    \texttt{Annotated[list, operator.add]}.
\end{itemize}

\section{Couche d'abstraction LLM}

\subsection{Interface \texttt{LLMClient}}

L'interface \texttt{LLMClient} est définie comme un Protocol Python
dans \texttt{src/llm/base.py}. Elle expose deux méthodes :

\begin{lstlisting}[caption={Interface LLMClient (src/llm/base.py).}]
class LLMClient(Protocol):
    def chat(
        self,
        system: str,
        messages: list[dict],
        temperature: float = 0.3,
        max_tokens: int = 2048,
    ) -> LLMResponse: ...

    def chat_structured(
        self,
        system: str,
        messages: list[dict],
        schema: type[T],
        temperature: float = 0.1,
    ) -> tuple[T, LLMResponse]: ...
\end{lstlisting}

\subsection{Quatre fournisseurs implémentés}

Quatre implémentations de \texttt{LLMClient} sont disponibles,
présentées dans le tableau~\ref{tab:providers}.

\begin{table}[htbp]
  \centering
  \caption{Fournisseurs LLM intégrés et méthode de structured output.}
  \label{tab:providers}
  \rowcolors{2}{lightgray}{white}
  \begin{tabular}{llll}
    \toprule
    \textbf{Provider} & \textbf{Modèles} & \textbf{Structured output} & \textbf{Coût relatif} \\
    \midrule
    Anthropic & claude-sonnet-4-6 & \texttt{tool\_use} forcé & $\times 10$ \\
    OpenAI & gpt-4o-mini & \texttt{response\_format} JSON & $\times 3$ \\
    Groq & llama-3.3-70b & \texttt{response\_format} JSON & $\times 1$ \\
    Nebius & DeepSeek-V3.2, Qwen3 & \texttt{response\_format} JSON & $\times 0.5$ \\
    \bottomrule
  \end{tabular}
\end{table}

La fabrique \texttt{get\_llm\_for\_task(task)} dans \texttt{src/llm/factory.py}
sélectionne automatiquement le modèle adapté à chaque type de tâche
selon la table \texttt{TASK\_MODELS} définie dans \texttt{src/config.py}.
Cette approche permet d'optimiser le rapport coût/qualité : Claude Sonnet
pour la synthèse (contexte 200K), Nebius Qwen3 pour la summarisation
répétitive (\$0.14/M tokens), GPT-4o-mini pour l'expansion de requêtes.

\section{Les sept agents}
\label{sec:agents}

\subsection{QueryExpander — Expansion de requêtes (ReAct)}

Le \texttt{QueryExpander} est le premier agent du pipeline. Il reçoit
le topic brut de l'utilisateur et génère plusieurs requêtes de recherche
diversifiées, optimisées pour maximiser la couverture des travaux pertinents.
Il implémente le pattern ReAct (détaillé au Chapitre~\ref{ch:patterns}).

À chaque itération, l'agent :
\begin{enumerate}
  \item Raisonne sur la meilleure requête à formuler ;
  \item Sonde OpenAlex avec cette requête ;
  \item Observe le nombre d'articles retournés et les concepts associés ;
  \item Adapte sa stratégie (affiner, élargir ou diversifier).
\end{enumerate}

À l'issue de 4 itérations maximum, le \texttt{QueryExpander} produit
une liste de 3 à 8 requêtes complémentaires qui alimentent le \texttt{Searcher}.

\subsection{Searcher — Récupération et déduplication}

Le \texttt{Searcher} exécute en parallèle toutes les requêtes produites
par le \texttt{QueryExpander} sur l'API OpenAlex, fusionne les résultats
et déduplique par identifiant OpenAlex. Il enrichit également les données
auteurs (h-index, affiliations) via des appels secondaires à l'API.

Le corpus résultant contient typiquement entre 80 et 300 articles, selon
la richesse du domaine et le nombre de requêtes générées.

\subsection{QualityCritic — Filtrage par qualité}

Le \texttt{QualityCritic} applique le scorer multi-dimensionnel
(Chapitre~\ref{ch:automl}) à chaque article du corpus. Il charge
automatiquement les pondérations Optuna apprises pour le domaine correspondant
(via \texttt{load\_weights\_for\_topic}), ou utilise les pondérations par
défaut si aucune pondération spécifique n'est disponible.

Il retient les $k$ articles les mieux classés (typiquement $k = 15$)
pour alimenter le \texttt{Summarizer}.

\subsection{Summarizer — Résumés structurés et narratifs}

Le \texttt{Summarizer} génère, pour chaque article sélectionné, un résumé
dans l'un des deux modes suivants :

\paragraph{Mode structuré.}
Le résumé est un objet \texttt{ArticleSummary} comportant six champs :
\textit{problem} (problème adressé), \textit{method} (approche proposée),
\textit{dataset} (données utilisées), \textit{results} (résultats obtenus),
\textit{limitations} (limites identifiées), et \textit{contribution}
(apport principal).

\paragraph{Mode narratif.}
Le résumé est un paragraphe en prose de 150 à 250 mots (\texttt{NarrativeSummary}),
plus accessible et adapté à une présentation à des lecteurs non spécialistes.

\subsection{Synthesizer — Synthèse globale (Reflexion)}

Le \texttt{Synthesizer} reçoit l'ensemble des résumés et génère une
synthèse globale du domaine (\texttt{Synthesis}) comprenant : une vue
d'ensemble du domaine, les principales approches méthodologiques,
les datasets de référence, et les résultats clés.

Il est au cœur du pattern Reflexion (Chapitre~\ref{ch:patterns}) :
sa production peut être révisée itérativement en réponse aux retours du
\texttt{Critic}.

\subsection{Critic — Évaluation Reflexion}

Le \texttt{Critic} évalue la synthèse produite par le \texttt{Synthesizer}
selon quatre axes : \textit{fidelity} (fidélité aux sources), \textit{completeness}
(exhaustivité), \textit{specificity} (niveau de détail), et \textit{consistency}
(cohérence interne). Il produit un objet \texttt{CriticFeedback} indiquant
si une révision est nécessaire (\texttt{needs\_revision}) et, le cas échéant,
les améliorations à apporter.

\subsection{TrendAnalyst — Analyse des tendances}

Le \texttt{TrendAnalyst} reçoit la synthèse finale et génère un objet
\texttt{TrendAnalysis} comprenant : les tendances émergentes du domaine
(avec leur niveau de maturité), les lacunes de recherche identifiées,
et des perspectives de directions futures. Un mécanisme de retry automatique
garantit que le champ \texttt{future\_perspectives} n'est jamais vide.

\section{Flux d'exécution et arêtes conditionnelles}

Le pipeline est compilé par LangGraph en un graphe orienté acyclique
(DAG) avec une arête conditionnelle pour le cycle Reflexion.
La figure~\ref{fig:langgraph_pipeline} illustre le flux complet.

\begin{figure}[htbp]
  \centering
  \includegraphics[width=0.85\textwidth]{figures/langgraph_pipeline.pdf}
  \caption{Pipeline LangGraph avec ses sept agents et l'arête conditionnelle
  Reflexion. L'arête \texttt{route\_after\_critic} redirige vers le
  \texttt{Synthesizer} si \texttt{needs\_revision = True} et que le
  nombre d'itérations est inférieur au maximum autorisé (3).}
  \label{fig:langgraph_pipeline}
\end{figure}

La fonction \texttt{route\_after\_critic} dans \texttt{src/agents/graph.py}
implémente la logique de routage :
\begin{itemize}
  \item Si \texttt{critic\_feedback.needs\_revision = True} et
    \texttt{reflexion\_count < MAX\_REFLEXION\_ITERATIONS} : retour au
    \texttt{Synthesizer} ;
  \item Sinon (qualité suffisante ou itérations épuisées) : passage au
    \texttt{TrendAnalyst}.
\end{itemize}

\section{Évaluation du pipeline}

Un exemple d'exécution du pipeline complet sur le topic
``\textit{fake news detection}'' produit les métriques suivantes
(tableau~\ref{tab:pipeline_demo}) :

\begin{table}[htbp]
  \centering
  \caption{Métriques d'exécution typiques du pipeline (topic: fake news detection).}
  \label{tab:pipeline_demo}
  \rowcolors{2}{lightgray}{white}
  \begin{tabular}{lrr}
    \toprule
    \textbf{Agent} & \textbf{Durée (s)} & \textbf{Coût (\$)} \\
    \midrule
    QueryExpander (4 iter.) & 12.3 & 0.002 \\
    Searcher (3 requêtes) & 8.7 & 0.000 \\
    QualityCritic (120 articles) & 0.3 & 0.000 \\
    Summarizer (15 articles) & 45.2 & 0.031 \\
    Synthesizer + Critic (2 iter.) & 18.9 & 0.012 \\
    TrendAnalyst & 7.4 & 0.005 \\
    \midrule
    \textbf{Total} & \textbf{92.8} & \textbf{0.050} \\
    \bottomrule
  \end{tabular}
\end{table}

Le coût total d'un run est de l'ordre de \$0.05 avec la configuration
Nebius/Groq, ce qui est compatible avec un usage régulier dans le cadre
d'un projet de recherche académique.

%% ============================================================
%%  Chapitre 6 — Patterns d'agents avancés
%% ============================================================
\chapter{Patterns d'agents avancés}
\label{ch:patterns}

Ce chapitre détaille les trois patterns de raisonnement avancés intégrés
dans \textit{Scientific Watch Agent} : le pattern ReAct pour l'expansion
de requêtes, le pattern Reflexion pour la révision de synthèse, et
ProTeGi pour l'optimisation automatique des prompts.

\section{Pattern ReAct : le QueryExpander}
\label{sec:react}

\subsection{Motivation}

Un utilisateur formule son besoin de recherche en termes naturels :
``\textit{fake news detection}''. Cette requête brute, soumise directement
à OpenAlex, peut manquer des travaux importants formulés différemment
(``\textit{misinformation}'' ou ``\textit{rumor detection}'').
L'expansion de requêtes consiste à générer automatiquement plusieurs
reformulations diversifiées du topic, de façon à maximiser la couverture
de la littérature pertinente.

Le pattern ReAct \cite{yao2022react} est particulièrement adapté à cette
tâche car il permet au LLM de \textit{tester} ses hypothèses en sondant
réellement la base de données avant de formuler la requête suivante,
évitant ainsi les expansions basées sur des présupposés erronés.

\subsection{Boucle ReAct dans le QueryExpander}

La boucle ReAct du \texttt{QueryExpander} fonctionne comme suit :

\begin{algorithm}[htbp]
  \caption{Boucle ReAct du QueryExpander}
  \label{alg:react}
  \begin{algorithmic}[1]
    \State $\text{queries} \gets [\text{topic}]$  \Comment{Requête initiale forcée}
    \State $i \gets 0$
    \While{$i < \text{MAX\_ITERATIONS}$ \textbf{and} $\text{stop\_reason} \neq \text{"done"}$}
      \State $\text{thought, query, stop} \gets \text{LLM}(\text{topic, history, observations})$
      \State $\text{count, concepts, examples} \gets \text{OpenAlex.probe}(\text{query})$
      \State $\text{observations} \mathrel{+}= (\text{query, count, concepts, examples})$
      \If{$\text{query} \notin \text{queries}$}
        \State $\text{queries} \mathrel{+}= [\text{query}]$
      \EndIf
      \State $i \mathrel{+}= 1$
    \EndWhile
    \State \Return $\text{queries}$
  \end{algorithmic}
\end{algorithm}

À chaque itération, l'agent LLM reçoit le contexte complet (topic initial,
historique des pensées, résultats des sondes) et produit une sortie
structurée \texttt{ReActThoughtAction} comprenant :
\begin{itemize}
  \item \texttt{thought} : son raisonnement sur la requête à formuler ;
  \item \texttt{search\_query} : la requête concrète pour OpenAlex ;
  \item \texttt{stop\_reason} : ``continue'' ou ``done'' (si l'exploration
    est jugée suffisante).
\end{itemize}

\paragraph{Contrainte sur la première itération.}
La requête de la première itération est forcée à être le topic original
exact, garantissant que le corpus inclura toujours les articles les plus
directement pertinents avant d'explorer des reformulations.

\subsection{Observation : le feedback de l'environnement}

L'observation renvoyée à l'agent après chaque sonde contient :
\begin{itemize}
  \item Le nombre d'articles retournés par OpenAlex ;
  \item Les cinq concepts les plus associés aux résultats ;
  \item Les titres des trois premiers articles.
\end{itemize}

Ces informations permettent à l'agent d'évaluer la qualité de sa requête
et d'ajuster sa stratégie. Par exemple, si une requête retourne moins de
10 articles, l'agent peut choisir d'élargir la formulation. Si les concepts
associés ne correspondent pas au domaine visé, il peut reformuler en
ajoutant des termes discriminants.

\section{Pattern Reflexion : le couple Synthesizer--Critic}
\label{sec:reflexion}

\subsection{Motivation}

La synthèse d'un corpus de publications est une tâche complexe, sujette
à des défauts difficiles à détecter automatiquement : omissions d'approches
importantes, généralisations abusives, formulations vagues, ou inconsistances
entre sections. Le pattern Reflexion \cite{shinn2023reflexion} adresse ce
problème en introduisant un agent critique qui évalue la synthèse et
formule des retours verbaux que le synthétiseur peut intégrer dans une
révision.

\subsection{Cycle Reflexion}

Le cycle Reflexion implique deux agents en boucle :

\begin{enumerate}
  \item Le \textbf{Synthesizer} génère ou révise la synthèse du domaine
    (\texttt{Synthesis}) en prenant en compte le retour du \texttt{Critic}
    du cycle précédent (si disponible) ;
  \item Le \textbf{Critic} évalue la synthèse reçue selon quatre axes
    et produit un retour structuré \texttt{CriticFeedback} ;
  \item Si \texttt{needs\_revision = True} et que le nombre d'itérations
    est inférieur à \texttt{MAX\_REFLEXION\_ITERATIONS = 3}, l'arête
    conditionnelle de LangGraph renvoie vers le \texttt{Synthesizer}.
\end{enumerate}

\subsection{Axes d'évaluation du Critic}

Le \texttt{Critic} évalue la synthèse sur quatre axes, chacun noté de
1 à 5 par le LLM :

\begin{description}
  \item[Fidelity (fidélité).] La synthèse est-elle fidèle au contenu
    des articles résumés ? Ne contient-elle pas d'affirmations non
    supportées par les sources ?
  \item[Completeness (exhaustivité).] La synthèse couvre-t-elle les
    principales approches et contributions du domaine, sans omission
    majeure ?
  \item[Specificity (spécificité).] La synthèse est-elle suffisamment
    précise, avec des noms de méthodes, des chiffres et des exemples
    concrets ? Évite-t-elle les généralités vagues ?
  \item[Consistency (cohérence).] Les différentes sections de la synthèse
    sont-elles cohérentes entre elles ? Y a-t-il des contradictions ?
\end{description}

Une révision est déclenchée si la qualité globale est ``poor'' (score
moyen $< 3$) ou ``acceptable'' (score moyen $< 4$).
La qualité ``good'' (score $\geq 4$) arrête le cycle.

\subsection{Mécanisme de sécurité}

En cas d'échec de l'appel LLM du \texttt{Critic} (timeout, erreur API),
le système approuve automatiquement la synthèse courante (\texttt{force\_approve}),
permettant au pipeline de continuer sans bloquer. Ce mécanisme garantit
la robustesse du système en production.

\subsection{Exemple d'itération Reflexion}

Le tableau~\ref{tab:reflexion_example} illustre l'évolution qualitative
d'une synthèse sur le domaine ``anomaly detection'' après une itération
Reflexion.

\begin{table}[htbp]
  \centering
  \caption{Exemple de retour Critic et amélioration après Reflexion.}
  \label{tab:reflexion_example}
  \begin{tabular}{p{0.45\textwidth}p{0.45\textwidth}}
    \toprule
    \textbf{Synthèse initiale (iter. 1)} & \textbf{Retour Critic} \\
    \midrule
    ``Les méthodes de détection d'anomalies utilisent des approches
    diverses incluant l'apprentissage profond et les méthodes statistiques.''
    &
    ``Specificity faible (2/5) : aucune méthode nommée.
    Completeness insuffisante (2/5) : autoencoders, Isolation Forest,
    VAE non mentionnés.'' \\
    \bottomrule
    \multicolumn{2}{l}{\textbf{Synthèse révisée (iter. 2)}} \\
    \multicolumn{2}{p{0.92\textwidth}}{``Les approches dominantes incluent
    les autoencoders et VAE (Variational Autoencoders) pour la détection
    non supervisée, Isolation Forest pour la détection rapide à haute dimension,
    et LSTN-AD pour les séries temporelles. Les benchmarks MVTec AD et
    KDD Cup 1999 sont les plus utilisés.''} \\
  \end{tabular}
\end{table}

\section{ProTeGi : optimisation automatique des prompts}
\label{sec:protegi}

\subsection{Motivation}

Le \texttt{Summarizer} utilise un prompt qui décrit la tâche de résumé
au LLM. La qualité de ce prompt affecte directement la qualité des résumés
produits, et donc la qualité de l'ensemble du pipeline. Une optimisation
manuelle est possible mais coûteuse et dépendante du modèle utilisé.
ProTeGi \cite{pryzant2023automatic} permet d'automatiser cette optimisation.

\subsection{Fonctionnement de ProTeGi}

ProTeGi adapte le concept de descente de gradient au domaine textuel.
L'algorithme procède par itérations :

\begin{enumerate}
  \item \textbf{Évaluation} : le prompt courant est utilisé pour résumer
    un batch d'articles ; les résumés sont évalués par rapport à des
    références humaines via ROUGE-L \cite{lin2004rouge} et BERTScore
    \cite{zhang2020bertscore} ;
  \item \textbf{Gradient textuel} : un LLM ``critique'' analyse les
    erreurs des résumés produits et formule des suggestions d'amélioration
    du prompt ;
  \item \textbf{Mise à jour} : les meilleures suggestions sont incorporées
    dans une nouvelle version du prompt ;
  \item \textbf{Sélection} : le meilleur prompt parmi les candidats est
    retenu pour l'itération suivante.
\end{enumerate}

\subsection{Résultats expérimentaux}

Nous avons conduit 14 runs d'optimisation ProTeGi sur 6 configurations
de LLM différentes (claude-haiku + claude-sonnet, gpt-4o-mini + gpt-4o,
llama-3.1-8b + gpt-4o, llama-3.1-8b + claude-sonnet, deepseek-v3.2,
nebius-qwen3), en modes structuré et narratif.

Le tableau~\ref{tab:protegi_results} présente les résultats obtenus
en mode narratif, qui s'est avéré le plus bénéficié par l'optimisation.

\begin{table}[htbp]
  \centering
  \caption{Résultats ProTeGi (mode narratif) — ROUGE-L et BERTScore.}
  \label{tab:protegi_results}
  \rowcolors{2}{lightgray}{white}
  \begin{tabular}{lccc}
    \toprule
    \textbf{Configuration} & \textbf{ROUGE-L init.} & \textbf{ROUGE-L final} & \textbf{$\Delta$} \\
    \midrule
    haiku + sonnet & 0.31 & 0.38 & $+22\,\%$ \\
    gpt-4o-mini + gpt-4o & 0.29 & 0.35 & $+21\,\%$ \\
    llama-8b + gpt-4o & 0.27 & 0.33 & $+22\,\%$ \\
    deepseek-v3.2 & 0.33 & 0.37 & $+12\,\%$ \\
    nebius-qwen3 & 0.30 & 0.36 & $+20\,\%$ \\
    \bottomrule
  \end{tabular}
\end{table}

L'amélioration moyenne en ROUGE-L est de $+19\,\%$ sur les 5 configurations
testées. Le gain est plus modeste en mode structuré ($+8\,\%$ en moyenne),
ce format étant plus contraint et donc moins sensible à la formulation
du prompt.

\paragraph{Analyse qualitative.}
L'examen des prompts optimisés révèle que ProTeGi a systématiquement
ajouté des instructions explicites sur : la longueur cible du résumé
(``\textit{entre 150 et 250 mots}''), la nécessité de mentionner les
performances numériques (``\textit{inclure les métriques quantitatives
reportées}''), et l'interdiction des phrases vagues (``\textit{éviter
les généralisations sans exemples concrets}'').

\section{Discussion}

Les trois patterns présentés dans ce chapitre jouent des rôles complémentaires
dans la qualité globale du système.

ReAct améliore la \textbf{couverture} en diversifiant les requêtes de recherche,
réduisant le risque de manquer des travaux importants formulés différemment.
Reflexion améliore la \textbf{précision} de la synthèse en permettant une
révision guidée par un retour structuré. ProTeGi améliore la \textbf{qualité
de présentation} des résumés en optimisant les instructions données au LLM.

Ces trois améliorations sont orthogonales et cumulatives : un système intégrant
les trois patterns bénéficie simultanément d'une meilleure couverture,
d'une meilleure précision et d'une meilleure qualité de présentation.

%% ============================================================
%%  Chapitre 8 — Méta-learning cross-domaine
%% ============================================================
\chapter{Méta-learning cross-domaine}
\label{ch:metalearning}

Ce chapitre présente la contribution scientifique centrale du mémoire :
le méta-learning pour la généralisation cross-domaine du scorer de qualité.
Nous formalisons le problème, décrivons les meta-features extraites,
présentons les méta-learners évalués, et discutons les résultats de la
validation de l'hypothèse H3.

\section{Motivation et formalisation}
\label{sec:meta_motivation}

\subsection{Le problème de la généralisation}

Le Chapitre~\ref{ch:automl} a démontré que les pondérations optimales
du scorer varient significativement d'un domaine à l'autre (tableau~\ref{tab:weight_profiles}).
Cette observation soulève un problème pratique : lorsqu'un utilisateur
soumet un nouveau domaine de recherche, il est impossible de lancer une
optimisation Optuna complète (qui nécessite un oracle préalablement construit)
avant d'avoir produit les premiers résultats.

Deux approches naïves existent pour contourner ce problème :
\begin{enumerate}
  \item \textbf{Default weights} : utiliser des pondérations uniformes,
    ignorant toute connaissance accumulée sur les autres domaines.
  \item \textbf{Transfer Direct} : utiliser la moyenne des pondérations
    apprises sur tous les domaines connus, supposant qu'une pondération
    ``moyenne'' sera approximativement correcte pour tout nouveau domaine.
\end{enumerate}

Ces deux approches sont des baselines. Nous proposons une troisième voie :
\textbf{le méta-learning}, qui prédit les pondérations optimales d'un nouveau
domaine à partir de caractéristiques mesurables du corpus de ce domaine,
sans nécessiter ni oracle ni optimisation préalable.

\subsection{Formalisation}

Soient $\mathcal{D} = \{d_1, \ldots, d_N\}$ les $N = 19$ domaines d'évaluation.
Pour chaque domaine $d_i$, nous disposons de :
\begin{itemize}
  \item Un vecteur de meta-features $\phi(d_i) \in \mathbb{R}^M$ décrivant
    les caractéristiques statistiques du corpus de $d_i$ ;
  \item Les pondérations optimales $\mathbf{w}^*(d_i) \in \mathcal{W}$
    apprises par Optuna sur l'oracle de $d_i$.
\end{itemize}

Le méta-learner est une fonction $f_\theta : \mathbb{R}^M \to \mathcal{W}$
apprise sur $\{(\phi(d_i), \mathbf{w}^*(d_i))\}_{i=1}^N$.
Pour un nouveau domaine $d_{\text{new}}$, les pondérations prédites sont
$\hat{\mathbf{w}}(d_{\text{new}}) = f_\theta(\phi(d_{\text{new}}))$.

\section{Extraction des meta-features}
\label{sec:meta_features}

\subsection{Définition des meta-features}

Les meta-features sont des statistiques descriptives calculées directement
sur le corpus de publications d'un domaine, \textbf{sans utiliser l'oracle}
(pour éviter toute fuite d'information). Elles capturent des propriétés
comme la maturité du domaine, la distribution des citations, ou le niveau
d'expertise des auteurs.

Le tableau~\ref{tab:meta_features_list} liste l'ensemble des 13 features
calculées dans notre système.

\begin{table}[htbp]
  \centering
  \caption{Meta-features extraites pour chaque domaine d'évaluation.}
  \label{tab:meta_features_list}
  \rowcolors{2}{lightgray}{white}
  \begin{tabular}{llp{0.45\textwidth}}
    \toprule
    \textbf{Feature} & \textbf{Type} & \textbf{Description} \\
    \midrule
    \texttt{grade2\_ratio} & ratio & Ratio d'articles grade-2 dans l'oracle \\
    \texttt{mean\_h\_index} & continu & H-index moyen des auteurs principaux \\
    \texttt{citation\_median} & continu & Médiane des citations par article \\
    \texttt{citation\_gini} & continu & Coefficient de Gini des citations \\
    \texttt{pct\_high\_hindex} & ratio & \% d'auteurs avec h-index $>$ 20 \\
    \texttt{pct\_recent} & ratio & \% d'articles des 3 dernières années \\
    \texttt{pct\_high\_cited} & ratio & \% d'articles avec $>$ 100 citations \\
    \texttt{year\_std} & continu & Écart-type des années de publication \\
    \texttt{gold\_ratio} & ratio & Ratio grade-2 / corpus total \\
    \texttt{median\_year} & entier & Année médiane de publication \\
    \texttt{unique\_author\_ratio} & ratio & Ratio d'auteurs uniques / total auteurs \\
    \texttt{pct\_q1} & ratio & \% d'articles publiés dans des revues Q1 \\
    \texttt{topic\_concept\_overlap} & ratio & Recouvrement topic / concepts OpenAlex \\
    \bottomrule
  \end{tabular}
\end{table}

\subsection{Sélection des meta-features pertinentes}
\label{sec:feature_selection}

La sélection des features pertinentes est effectuée par analyse des
corrélations de Pearson entre chaque meta-feature et chaque dimension
de pondération, sur l'ensemble des $N = 19$ domaines. Le tableau~\ref{tab:pearson_correlations}
présente les corrélations maximales (en valeur absolue) obtenues.

\begin{table}[htbp]
  \centering
  \caption{Corrélations de Pearson (max $|r|$ sur les 6 pondérations), n=19 domaines.}
  \label{tab:pearson_correlations}
  \rowcolors{2}{lightgray}{white}
  \begin{tabular}{lccl}
    \toprule
    \textbf{Feature} & \textbf{max $|r|$} & \textbf{Sélectionnée ?}
      & \textbf{Corrélation principale} \\
    \midrule
    \texttt{grade2\_ratio}    & \textbf{0.67} & \checkmark & velocity ($r = -0.67$) \\
    \texttt{mean\_h\_index}   & \textbf{0.61} & \checkmark & relevance ($r = -0.61$) \\
    \texttt{citation\_median} & \textbf{0.58} & \checkmark & recency ($r = +0.58$) \\
    \texttt{citation\_gini}   & \textbf{0.56} & \checkmark & venue ($r = -0.56$) \\
    \texttt{pct\_high\_hindex}& \textbf{0.52} & \checkmark & relevance ($r = -0.52$) \\
    \texttt{pct\_recent}      & \textbf{0.45} & \checkmark & relevance ($r = -0.45$) \\
    \texttt{pct\_high\_cited} & \textbf{0.41} & \checkmark & recency ($r = +0.41$) \\
    \texttt{gold\_ratio}      & 0.39 & \texttimes & — \\
    \texttt{year\_std}        & 0.38 & \texttimes & — \\
    \texttt{unique\_author\_ratio} & 0.18 & \texttimes & — \\
    \bottomrule
  \end{tabular}
\end{table}

Les 7 features avec $\max|r| > 0.40$ constituent l'\texttt{EFFECTIVE\_FEATURES}
utilisé par le méta-learner, défini dans \texttt{src/metalearning/meta\_learner.py}.

\paragraph{Interprétation des corrélations.}
La corrélation la plus forte est celle entre \texttt{grade2\_ratio} et la
pondération velocity ($r = -0.67$) : les domaines avec une plus grande
proportion d'articles de référence (\textit{landmarks}) accordent moins
d'importance à la vélocité des citations, car les articles fondateurs restent
pertinents même s'ils ne sont plus activement cités récemment. La corrélation
entre \texttt{mean\_h\_index} et la pondération relevance ($r = -0.61$) indique
que les domaines à experts très reconnus (h-index élevé) ont moins besoin
de la pertinence sémantique : la réputation de l'auteur est suffisamment
discriminante.

\section{Méta-learners évalués}
\label{sec:meta_learners}

Nous évaluons 12 méta-learners, répartis en quatre familles.

\subsection{Méthodes à base de voisins (kNN)}

Le méta-learner kNN($k$) prédit les pondérations d'un domaine test $d$
comme une moyenne pondérée par la distance inverse des pondérations des
$k$ domaines d'entraînement les plus proches dans l'espace des meta-features :

\begin{equation}
  \hat{\mathbf{w}}_{\text{kNN}(k)}(d) \propto
  \sum_{i \in \mathcal{N}_k(d)} \frac{1}{\|\phi(d) - \phi(d_i)\|_2}
  \cdot \mathbf{w}^*(d_i)
  \label{eq:knn_prediction}
\end{equation}

où $\mathcal{N}_k(d)$ est l'ensemble des $k$ voisins les plus proches.
Les meta-features sont normalisées par MinMaxScaler avant le calcul de
distance, pour éviter que \texttt{citation\_median} (grande amplitude)
ne domine l'espace métrique.

Nous évaluons $k \in \{1, 2, 3\}$.

\subsection{Méthodes de régression linéaire}

Trois régresseurs linéaires sont évalués, chacun ajusté indépendamment
pour chaque dimension de pondération :

\paragraph{RidgeCV.} Régression ridge avec validation croisée automatique
pour la sélection du paramètre de régularisation $\alpha$ parmi
$\{0.01, 0.1, 1, 10, 100, 1000\}$.

\paragraph{BayesianRidge.} Régression bayésienne de crête qui ajuste
automatiquement les hyperparamètres de régularisation via l'inférence
bayésienne (maximisation de la log-vraisemblance marginale).

\paragraph{GPR (\textit{Gaussian Process Regression}).}
Régression par processus gaussien avec noyau RBF + constant + bruit blanc.
Cette méthode modélise l'incertitude de prédiction et est théoriquement
bien adaptée aux données de faible taille ($n = 18$ en LOO), mais souffre
en pratique d'instabilité numérique sur certains domaines.

\subsection{Méthodes d'ensemble (rank fusion)}

Ces méthodes combinent les prédictions des quatre base-learners
kNN(1), kNN(2), kNN(3) et BayesianRidge en un vecteur de pondérations
unique.

\paragraph{Soft Vote.} Moyenne arithmétique des prédictions des 4 membres,
suivie d'une normalisation :
\begin{equation}
  \hat{\mathbf{w}}_{\text{SoftVote}} = \frac{1}{4} \sum_{m=1}^{4} \hat{\mathbf{w}}_m
  \label{eq:soft_vote}
\end{equation}

\paragraph{Hard Vote (médiane).} Médiane par dimension sur les 4 membres
(avec 4 valeurs : moyenne des 2e et 3e valeurs ordonnées), suivie d'une
normalisation. Plus robuste aux prédictions aberrantes (\textit{outliers})
qu'une moyenne.

\paragraph{Borda Count.} Pour chaque membre $m$, les 6 dimensions de
pondération sont classées par valeur décroissante. La dimension en
position $r$ (sur 6) reçoit un score de Borda $b = 6 - 1 - r = 5 - r$.
Les scores sont sommés sur les 4 membres et normalisés :

\begin{equation}
  \hat{\mathbf{w}}_{\text{Borda}} \propto \sum_{m=1}^{4}
  \mathbf{b}(\hat{\mathbf{w}}_m)
  \label{eq:borda}
\end{equation}

où $\mathbf{b}(\mathbf{w})_j = 5 - \text{rank}_j(\mathbf{w})$.

\paragraph{Reciprocal Rank Fusion (RRF).} Score RRF par dimension :
\begin{equation}
  \text{RRF}_j = \sum_{m=1}^{4} \frac{1}{k + \text{rank}_j(\hat{\mathbf{w}}_m)}
  \label{eq:rrf}
\end{equation}
avec $k = 1$ (choix adapté à $n = 6$ dimensions : $k = 60$ standard
sur-lisserait les différences).

\subsection{Ensembles hybrides kNN--régression}

Deux ensembles convexes combinent kNN(1) avec un régresseur linéaire :

\begin{align}
  \hat{\mathbf{w}}_{\text{En+R}}(\alpha) &= \alpha \cdot \hat{\mathbf{w}}_{\text{kNN}(1)}
    + (1 - \alpha) \cdot \hat{\mathbf{w}}_{\text{Ridge}} \\
  \hat{\mathbf{w}}_{\text{En+B}}(\alpha) &= \alpha \cdot \hat{\mathbf{w}}_{\text{kNN}(1)}
    + (1 - \alpha) \cdot \hat{\mathbf{w}}_{\text{Bayes}}
\end{align}

Le paramètre $\alpha \in \{0.0, 0.1, \ldots, 1.0\}$ est sélectionné par
balayage sur le jeu d'entraînement LOO.

\section{Protocole d'évaluation : LOO-CV}

Chaque méta-learner est évalué par validation croisée Leave-One-Out sur
les 19 domaines. À chaque fold $i$ :
\begin{enumerate}
  \item Le méta-learner est ajusté sur les 18 domaines $\mathcal{D} \setminus \{d_i\}$ ;
  \item Il prédit $\hat{\mathbf{w}}(d_i)$ à partir de $\phi(d_i)$ ;
  \item Le NDCG@15 est calculé en classant les articles du corpus $C_{d_i}$
    avec les pondérations prédites, en utilisant les sous-scores précalculés
    pour accélérer le calcul.
\end{enumerate}

La performance d'un méta-learner sur un domaine est comparée au Transfer
Direct (pondérations moyennes sur les 18 domaines d'entraînement).
Un méta-learner remporte une ``victoire'' sur un domaine s'il surpasse
le Transfer Direct de plus de $10^{-6}$ en NDCG@15.

\section{Résultats : validation de H3}
\label{sec:h3_results}

\subsection{Taux de victoire par méthode}

Le tableau~\ref{tab:loo_results} présente les taux de victoire (vs Transfer
Direct) et les moyennes de NDCG@15 pour les 12 méthodes évaluées.

\begin{table}[htbp]
  \centering
  \caption{Résultats LOO-CV : victoires vs Transfer Direct et NDCG@15 moyen.
  Les méthodes dépassant le seuil de 75\,\% sont mises en évidence.}
  \label{tab:loo_results}
  \rowcolors{2}{lightgray}{white}
  \begin{tabular}{lccr}
    \toprule
    \textbf{Méthode} & \textbf{Wins/19} & \textbf{\%} & \textbf{Mean NDCG@15} \\
    \midrule
    \rowcolor{green!20}
    \textbf{Borda} & \textbf{15/19} & \textbf{79\,\%} $\checkmark$ & \textbf{0.282} \\
    RRF & 14/19 & 74\,\% & 0.254 \\
    En+R ($\alpha = 0.1$) & 13/19 & 68\,\% & 0.253 \\
    BayesRidge & 12/19 & 63\,\% & 0.253 \\
    En+B ($\alpha = 0.6$) & 12/19 & 63\,\% & 0.258 \\
    kNN(2) & 11/19 & 58\,\% & 0.249 \\
    kNN(3) & 11/19 & 58\,\% & 0.254 \\
    Ridge & 11/19 & 58\,\% & 0.249 \\
    SoftVote & 10/19 & 53\,\% & 0.252 \\
    HardVote & 9/19 & 47\,\% & 0.235 \\
    kNN(1) & 8/19 & 42\,\% & 0.259 \\
    GPR & 7/19 & 37\,\% & 0.235 \\
    \midrule
    \textit{Transfer Direct} & \textit{(référence)} & — & 0.247 \\
    \textit{Default weights} & \textit{(référence)} & — & 0.229 \\
    \textit{Optuna*} & \textit{(borne sup.)} & — & 0.387 \\
    \bottomrule
  \end{tabular}
  \small{*Optuna = pondérations réelles apprises sur le domaine test (non disponibles en pratique).}
\end{table}

\begin{resultbox}
  \textbf{H3 confirmée.} Borda Count atteint 15/19 = 79\,\% de taux de
  victoire sur le Transfer Direct, dépassant le seuil de 75\,\% fixé par H3.
  Sa moyenne de NDCG@15 est 0.282, soit $+14.2\,\%$ par rapport au Transfer
  Direct (0.247) et $+23.1\,\%$ par rapport aux pondérations par défaut (0.229).
\end{resultbox}

\subsection{Analyse par domaine}

La figure~\ref{fig:wins_heatmap} présente le tableau détaillé des gains
($\Delta$NDCG@15 vs Transfer Direct) par domaine et par méthode.
Les domaines sont triés par delta Borda décroissant.

\begin{figure}[htbp]
  \centering
  \includegraphics[width=\textwidth]{figures/diag_method_wins.png}
  \caption{Tableau des gains NDCG@15 (vs Transfer Direct) par domaine et méthode.
  Vert = victoire, rouge = défaite, ligne séparatrice après le 15\textsuperscript{e}
  domaine (médiane Borda).}
  \label{fig:wins_heatmap}
\end{figure}

Quatre domaines résistent à l'ensemble des méta-learners :
\textit{Transfer Learning} ($\Delta_{\text{Borda}} = -0.009$),
\textit{Fake News Detection} ($\Delta_{\text{Borda}} = -0.029$),
\textit{Speech ASR} ($\Delta_{\text{Borda}} = -0.107$), et
\textit{Anomaly Detection} ($\Delta_{\text{Borda}} = -0.148$).
Ces domaines présentent des profiles atypiques que les meta-features
actuelles ne capturent pas fidèlement, ce qui est discuté à la
section~\ref{sec:meta_discussion}.

\section{Test statistique de Wilcoxon}
\label{sec:wilcoxon}

\subsection{Configuration du test}

Afin d'évaluer si l'amélioration de Borda sur le Transfer Direct est
statistiquement significative et non due au hasard, nous appliquons
le test de Wilcoxon signé \cite{wilcoxon1945individual} (voir
Section~\ref{sec:etat_art_wilcoxon} pour la définition formelle) :

\begin{itemize}
  \item \textbf{Observations appariées} : $\Delta_i = \text{NDCG}_{Borda}(d_i) -
    \text{NDCG}_{Transfer}(d_i)$ pour $i = 1, \ldots, 19$ ;
  \item \textbf{Test unilatéral} : $H_1 : \text{median}(\Delta) > 0$ ;
  \item \textbf{Seuil de significativité} : $\alpha = 0.05$.
\end{itemize}

\subsection{Résultats}

\begin{table}[htbp]
  \centering
  \caption{Test de Wilcoxon signé (unilatéral) : Borda vs Transfer Direct, $n=19$.}
  \label{tab:wilcoxon_results}
  \begin{tabular}{lrrrll}
    \toprule
    \textbf{Méthode} & \textbf{$W^+$} & \textbf{$p$-value} & \textbf{$r_{rb}$}
      & \textbf{Effet} & \textbf{Sig. ($\alpha=0.05$) ?} \\
    \midrule
    \textbf{Borda} & \textbf{149.0} & \textbf{0.0145} & \textbf{+0.568}
      & \textbf{large} & \textbf{\checkmark OUI} \\
    En+R & 130.5 & 0.077 & +0.368 & medium & non \\
    RRF & 131.0 & 0.078 & +0.379 & medium & non \\
    BayesRidge & 127.5 & 0.095 & +0.345 & medium & non \\
    En+B & 126.0 & 0.113 & +0.326 & medium & non \\
    Ridge & 117.0 & 0.198 & +0.232 & small & non \\
    kNN(1) & 96.0 & 0.492 & +0.011 & small & non \\
    GPR & 87.5 & 0.619 & $-$0.143 & small & non \\
    \bottomrule
  \end{tabular}
\end{table}

\begin{resultbox}
  Borda Count est la seule méthode statistiquement significative à $\alpha = 0.05$
  ($W^+ = 149$, $p = 0.014$, $r_{rb} = +0.568$, effet large).
  En+R, RRF et BayesRidge montrent des effets medium ($r_{rb}$ entre 0.33 et 0.38)
  et des $p$-values entre 0.07 et 0.11, suggérant des tendances positives
  qui ne peuvent être confirmées avec $n = 19$.
\end{resultbox}

La figure~\ref{fig:wilcoxon_plot} illustre les $p$-values et tailles d'effet
pour les 12 méthodes.

\begin{figure}[htbp]
  \centering
  \includegraphics[width=\textwidth]{figures/phase8d_wilcoxon.png}
  \caption{$p$-values (haut, barres vertes = significatives) et corrélations
  rang-bisériales $r_{rb}$ (bas, effet size) pour les 12 méta-learners.
  La ligne orange tiretée indique $\alpha = 0.05$.}
  \label{fig:wilcoxon_plot}
\end{figure}

\section{Discussion}
\label{sec:meta_discussion}

\subsection{Pourquoi Borda excelle-t-il ?}

Le succès du Borda Count s'explique par sa robustesse aux prédictions
aberrantes. En agrégeant des rangs plutôt que des valeurs absolues,
il réduit l'impact des cas où un membre du comité (kNN ou BayesRidge)
produit une prédiction très éloignée de la réalité. Pour $n = 18$ domaines
d'entraînement, la variabilité des prédictions individuelles est élevée,
et les méthodes d'agrégation de rangs sont connues pour être plus robustes
en régime de faible données \cite{dwork2001rank}.

\subsection{Domaines difficiles}

Les quatre domaines où Borda échoue (\textit{Transfer Learning}, \textit{Fake
News}, \textit{Speech ASR}, \textit{Anomaly Detection}) partagent une
caractéristique commune : leurs profils de pondération optimaux sont
atypiques par rapport au reste du corpus. Par exemple, \textit{Anomaly
Detection} requiert des pondérations très concentrées sur l'impact
($w_{\text{impact}} = 0.39$) et l'auteur ($w_{\text{authors}} = 0.28$),
un profil non prédit par les meta-features actuelles.

Ces cas difficiles suggèrent que l'ensemble actuel de meta-features
n'est pas suffisamment expressif pour capturer toute la diversité des
domaines. Des meta-features supplémentaires — comme le ratio de preprints
arXiv, la densité sémantique du corpus, ou l'entropie des concepts
OpenAlex — pourraient améliorer la couverture des cas extrêmes.

\subsection{Limites et perspectives}

\paragraph{Taille du corpus d'évaluation.}
Avec $n = 19$ domaines, la puissance statistique du test de Wilcoxon est
limitée. Une étude avec $n \geq 30$ domaines permettrait de confirmer la
significativité d'En+R et RRF, actuellement à la frontière du seuil ($p \approx 0.08$).

\paragraph{Qualité des oracles.}
La construction des oracles par heuristique bibliographique (citations
dans des surveys) introduit un biais : les articles non cités dans les
surveys disponibles, même s'ils sont pertinents, sont classés grade-0.
Une annotation humaine sur un sous-ensemble permettrait de valider la
qualité de ces oracles.

\paragraph{Meta-features dynamiques.}
Les meta-features actuelles sont calculées une seule fois sur le corpus
initial. Des meta-features dynamiques, recalculées après chaque requête
de l'utilisateur, pourraient mieux capturer les propriétés spécifiques
de la requête plutôt que du domaine général.

\paragraph{Extension au méta-learning de niveau 2.}
La contribution actuelle correspond au Niveau 3 de généralisation
(méta-apprentissage). Une extension vers le Niveau 2 (fine-tuning rapide
sur quelques exemples du nouveau domaine) permettrait d'affiner les
prédictions lorsqu'un minimum de feedback utilisateur est disponible.


%% ── Bibliographie ──────────────────────────────────────────
\printbibliography[heading=bibintoc, title={Références bibliographiques}]

%% ── Annexes ────────────────────────────────────────────────
\appendix
\chapter{Configuration des domaines d'évaluation}
\label{app:domains}

Le tableau~\ref{tab:domains-config} liste les 19 domaines utilisés dans l'évaluation
du méta-learner, avec le topic OpenAlex et l'année de début de la fenêtre temporelle.

\begin{longtable}{llc}
\caption{Configuration des 19 domaines d'évaluation.}
\label{tab:domains-config} \\
\toprule
\textbf{Domaine} & \textbf{Topic de recherche} & \textbf{Depuis} \\
\midrule
\endfirsthead
\multicolumn{3}{c}{\textit{(suite du tableau~\ref{tab:domains-config})}} \\
\toprule
\textbf{Domaine} & \textbf{Topic de recherche} & \textbf{Depuis} \\
\midrule
\endhead
\midrule
\multicolumn{3}{r}{\textit{Suite page suivante\ldots}} \\
\endfoot
\bottomrule
\endlastfoot
fake\_news\_detection & fake news detection misinformation & 2017 \\
graph\_neural\_networks & graph neural networks node classification & 2017 \\
federated\_learning & federated learning privacy distributed & 2017 \\
medical\_image\_segmentation & medical image segmentation deep learning & 2015 \\
neural\_architecture\_search & neural architecture search AutoML & 2017 \\
explainable\_ai & explainable AI interpretability SHAP LIME & 2017 \\
knowledge\_graph\_embedding & knowledge graph embedding link prediction & 2017 \\
graph\_attention\_networks & graph attention networks & 2018 \\
object\_detection\_deep\_learning & object detection YOLO deep learning & 2016 \\
sentiment\_analysis & sentiment analysis opinion mining BERT & 2018 \\
transfer\_learning & transfer learning domain adaptation & 2017 \\
generative\_adversarial\_networks & generative adversarial networks image synthesis & 2017 \\
anomaly\_detection & anomaly detection unsupervised & 2017 \\
deep\_reinforcement\_learning & deep reinforcement learning policy gradient & 2017 \\
text\_summarization & text summarization abstractive extractive & 2018 \\
speech\_recognition\_asr & speech recognition automatic ASR & 2017 \\
computational\_biology & machine learning genomics bioinformatics & 2017 \\
medical\_nlp & NLP clinical text biomedical & 2017 \\
quantitative\_finance\_ml & machine learning financial prediction & 2018 \\
\end{longtable}

\end{document}
